\documentclass[11pt,a4paper]{article}

%======================================================================
%  Quantum Computing for Data Science -- DU Mathematics, UCA 2026-2027
%  Corrected and completed version of the lecture notes of C. Kriegler
%======================================================================

\usepackage[utf8]{inputenc}
\usepackage[T1]{fontenc}
\usepackage{lmodern}
\usepackage[a4paper,margin=2.5cm]{geometry}
\usepackage{amsmath,amssymb,amsthm,mathtools}
\usepackage{enumitem}
\usepackage{longtable,array,booktabs}
\usepackage[dvipsnames]{xcolor}
\usepackage[most]{tcolorbox}
\usepackage{hyperref}
\hypersetup{colorlinks=true,linkcolor=BrickRed,citecolor=BrickRed,urlcolor=NavyBlue}

\numberwithin{equation}{section}
\setlist[enumerate]{label=\arabic*.,leftmargin=2em}

%---------------------------------------------------------------- macros
\newcommand{\N}{\mathbb{N}}
\newcommand{\Z}{\mathbb{Z}}
\newcommand{\R}{\mathbb{R}}
\newcommand{\C}{\mathbb{C}}
\newcommand{\K}{\mathbb{K}}
\newcommand{\PP}{\mathbb{P}}
\newcommand{\E}{\mathbb{E}}
\newcommand{\ip}[2]{\langle #1, #2\rangle}
\newcommand{\norm}[1]{\left\| #1\right\|}
\newcommand{\abs}[1]{\left| #1\right|}
\newcommand{\ket}[1]{|#1\rangle}
\newcommand{\bra}[1]{\langle #1|}
\newcommand{\ketbra}[2]{|#1\rangle\langle #2|}
\newcommand{\conj}[1]{\overline{#1}}
\DeclareMathOperator{\tr}{tr}
\DeclareMathOperator{\spn}{span}
\DeclareMathOperator{\rank}{rank}
\DeclareMathOperator{\Ran}{Ran}
\DeclareMathOperator{\Ker}{Ker}
\DeclareMathOperator{\Real}{Re}
\DeclareMathOperator{\Imag}{Im}
\newcommand{\one}{1}

%------------------------------------------------- theorem environments
% Same numbering scheme as the original notes (shared counter per section)
\theoremstyle{plain}
\newtheorem{theorem}{Theorem}[section]
\newtheorem{proposition}[theorem]{Proposition}
\newtheorem{lemma}[theorem]{Lemma}
\theoremstyle{definition}
\newtheorem{definition}[theorem]{Definition}
\newtheorem{example}[theorem]{Example}
\newtheorem{exercise}[theorem]{Exercise}
\newtheorem{remark}[theorem]{Remark}

% Results that are NOT in the original notes get a letter suffix
% (e.g. Lemma 1.7a) so that the original numbering is preserved.
\newcommand{\thistheoremname}{}
\theoremstyle{plain}
\newtheorem*{addedthm*}{\thistheoremname}
\theoremstyle{definition}
\newtheorem*{addeddef*}{\thistheoremname}
\newenvironment{addedthm}[1]{\renewcommand{\thistheoremname}{#1}\begin{addedthm*}}{\end{addedthm*}}
\newenvironment{addeddef}[1]{\renewcommand{\thistheoremname}{#1}\begin{addeddef*}}{\end{addeddef*}}

\renewcommand{\qedsymbol}{$\blacksquare$}
\newenvironment{solution}{\begin{proof}[Solution]}{\end{proof}}

%------------------------------------------------------- coloured boxes
\newtcolorbox{correction}[1][]{enhanced,breakable,colback=red!4,colframe=BrickRed,
  fonttitle=\bfseries\small,title={Correction #1},left=4pt,right=4pt,top=2pt,bottom=2pt}
\newtcolorbox{addition}[1][]{enhanced,breakable,colback=blue!3,colframe=NavyBlue!70,
  fonttitle=\bfseries\small,title={Complement #1},left=4pt,right=4pt,top=2pt,bottom=2pt}

\begin{document}

%======================================================================
\begin{titlepage}
\centering
\vspace*{3cm}
{\Large\bfseries Universit\'e Clermont Auvergne\\[2pt]
Quantum Computing for Data Science DU\\[2pt] Mathematics\\[2pt] Year 2026 -- 2027\par}
\vspace{1cm}
{\large Lecture notes of Christoph Kriegler (November 2026)\par}
\vspace{0.6cm}
{\large\bfseries Corrected and completed edition\par}
\vspace{0.3cm}
{\normalsize with all proofs, all exercise solutions and an errata list\par}
\vfill
\begin{minipage}{0.85\textwidth}\small
\textbf{How to read this edition.} The text of the original notes is reproduced with
the \emph{same numbering} (Definition 1.1, \dots, Example 8.6). In addition:
\begin{itemize}[leftmargin=1.5em]
\item \textcolor{BrickRed}{\textbf{red boxes}} (``Correction'') point out and fix a
  mistake of the original notes; the statement printed in the main text is already the
  corrected one;
\item \textcolor{NavyBlue}{\textbf{blue boxes}} (``Complement'') contain material that was
  missing or only announced (``one can prove'', ``with a little more work'', \dots);
\item every exercise is followed by a complete \emph{Solution};
\item every statement given without proof is followed by a \emph{Proof};
\item new auxiliary results carry a letter suffix (e.g.\ Lemma 1.7a) so that the
  original numbering is unchanged.
\end{itemize}
A summary of all corrections is given on the next pages.

\medskip
\textbf{Conventions.} $\N=\{1,2,3,\dots\}$ and $\N_0=\N\cup\{0\}$ (this is the convention
implicitly used in Example 3.5). $\K$ denotes $\R$ or $\C$. Scalar products are
\emph{anti-linear in the first} and linear in the second argument (physicists' convention,
as in Definition 1.6).
\end{minipage}
\vspace{1cm}
\end{titlepage}

\tableofcontents
\newpage
%======================================================================
\section*{Summary of corrections (errata)}
\addcontentsline{toc}{section}{Summary of corrections (errata)}

Severity: \textbf{M} = mathematical error (wrong formula/statement),
\textbf{G} = gap (true statement, incomplete argument or imprecise definition),
\textbf{T} = typo / wording.

{\small
\renewcommand{\arraystretch}{1.25}
\setlength{\tabcolsep}{4pt}
\begin{longtable}{@{}p{0.6cm}p{1.9cm}p{5.6cm}p{6.0cm}c@{}}
\toprule
\# & Where & Original & Corrected & \\
\midrule\endhead
1 & Def.~1.1 & The list of vector space axioms omits commutativity and associativity of $+$ and the distributive law $\lambda(x+y)=\lambda x+\lambda y$. & Axioms completed. & G\\
2 & Def.~1.6(3) & ``anti-symetry'' & \emph{conjugate (Hermitian) symmetry}; ``anti-symmetry'' means something else ($\ip{x}{y}=-\ip{y}{x}$). & T\\
3 & Ex.~1.12 & $\ip{x}{y}=\sum_{k=1}^n\conj{x_n}\,y_n$ & $\ip{x}{y}=\sum_{k=1}^n\conj{x_k}\,y_k$ (wrong summation index). & M\\
4 & after Ex.~1.14 & ``$C([a,b])$ is not a Hilbert space'' -- no justification. & Explicit non-convergent Cauchy sequence given. & G\\
5 & Rem.~1.21, Ex.~1.22, \dots & Bra/kets typeset with the relation symbols $<,>$ ($|\phi>$, $<\psi|$). & $\ket{\phi}$, $\bra{\psi}$ with angle brackets $\langle\ \rangle$. & T\\
6 & Prop.~2.8 & For a non-separable Hilbert space the orthonormal basis is uncountable; ``convergent series'' must be read as an unconditional sum with countably many non-zero terms. & Remark added. & G\\
7 & Ex.~3.4 & ``$X$ is called a \emph{Bernoulli} variable with parameters $n$ and $p$'' & This is a \emph{binomial} variable $\mathcal B(n,p)$. A Bernoulli variable is the case $n=1$. & M\\
8 & Ex.~3.6 & ``Show that $\PP(\{X\in\{0,\dots,n\}\})=1$'' & Trivially true since $X$ takes values in $\{0,\dots,n\}$; the intended (non-trivial) statement is $\sum_{k}\PP(A_k)=1$, i.e.\ the definition is consistent. Same for Poisson. & G\\
9 & Def.~3.11 & $\E(X)$ defined ``whenever the series converges'' & Must require \emph{absolute} convergence (a sum over a countable set has no natural order). & M\\
10 & Ex.~3.12(3) & Poisson parameter $\lambda\geqslant 0$, while Ex.~3.5 requires $\lambda>0$. & Harmless ($\lambda=0$ gives $X=0$ a.s.); made consistent. & T\\
11 & Prop.~4.1(2) & $a_{ij}=\ip{T(e_j)}{f_i}$ & $a_{ij}=\ip{f_i}{T(e_j)}$. With a scalar product anti-linear in the first variable, the original formula gives $\conj{a_{ij}}$. & M\\
12 & Prop.~4.1, Prop.~6.3, Ex.~4.7 & ``lines'' of a matrix; ``bijective (i.e.\ exact)''; ``quantum calculus'' & ``rows''; ``bijective''; ``quantum mechanics''. & T\\
13 & \S5 intro & ``$\lambda\in C$'' & $\lambda\in\C$. & T\\
14 & Def.~5.4 & $e^{tA}$ defined only for $t\geqslant 0$ & The proof of Prop.~5.7 uses $e^{hA}$ with $h<0$: define $e^{tA}$ for all $t\in\R$ (even $t\in\C$). & G\\
15 & Lemma 5.6 & Proof omitted. & Proof via Cauchy product given. & G\\
16 & Proof of Prop.~5.7 & ``according to Lemma ??'' & Broken reference: Lemma~5.6. & T\\
17 & Proof of Prop.~5.7 & $\big\|\sum_{n\ge1}A\frac{(hA)^{n-1}}{n!}u_1-Au_1\big\| = \big\|\sum_{n\ge2}\frac{(hA)^{n-1}}{n!}u_1\big\|$, final bound $|h|\,\|A\|\,e^{|h|\|A\|}\|u_1\|$ & A factor $A$ is lost: correct is $\big\|\sum_{n\ge2}A\frac{(hA)^{n-1}}{n!}u_1\big\|\le |h|\,\|A\|^2e^{|h|\|A\|}\|u_1\|$. Conclusion unchanged. & M\\
18 & Prop.~5.7 & Uniqueness ``with a little more work'' & Uniqueness proved. & G\\
19 & Ex.~5.8 & ``expoential'' (twice) & ``exponential''. & T\\
20 & Thm.~5.9 & Proof uses $C^k=e^{A+B}$ without saying so; uses $k$ instead of $n$. & Made explicit. & G\\
21 & Proof of Prop.~6.3 & ``It remains to show that $A$ is positive \emph{definite}'' & positive \emph{semi}-definite (e.g.\ a pure state is never positive definite for $n\ge2$). & M\\
22 & Proof of Prop.~6.3 & ``the matrix of $T_x$ is $(\conj{x_i}x_j)_{ij}$'' & The matrix of $T_x=\ketbra{x}{x}$ is $(x_i\conj{x_j})_{ij}$; then $\tr(AT_x)=\sum_{i,j}\conj{x_i}a_{ij}x_j=\ip{x}{Ax}$. & M\\
23 & Prop.~6.3 & Uniqueness and the converse are not proved. & Proved. & G\\
24 & Ex.~6.4 & $\sum_{k=1}^{n}\lambda_kT_{x_k}$ & $\sum_{k=1}^{N}\lambda_kT_{x_k}$. & T\\
25 & Ex.~6.5(1) & $A=\begin{psmallmatrix}2&-1&-1\\-1&5&1\\-1&1&1\end{psmallmatrix}$ is claimed to equal $T_x+T_y$ & False: $T_x+T_y=\begin{psmallmatrix}2&1&-1\\1&5&1\\-1&1&1\end{psmallmatrix}$ (sign of $a_{12}=a_{21}$). In fact no $x=(x_1,x_2,0)$, $y\in\R^3$ give the printed matrix. & M\\
26 & Ex.~6.5(2) & ``Deduce that $A$ is positive definite'' & False for $A=T_x+T_y$: it has rank 2 ($\det A=0$, kernel $\spn\{(-2,1,-3)\}$). It is positive \emph{semi}-definite. & M\\
27 & Ex.~6.5(2) & ``$\Phi:M_3(\C)\to M_3(\C)$, $B\mapsto\tr(AB)$ defines a mixed state'' & Codomain is $\C$, and $\tr A=8\neq1$: the state is $\Phi(B)=\frac18\tr(AB)$, density matrix $\rho=\frac18A$. & M\\
28 & Ex.~7.5 & ``$x,y$ are colinear'' & Means: $x,y$ and $0$ lie on one line (i.e.\ $x,y$ linearly dependent). ``open circles'' $\to$ ``open discs''. Triangle inequality not proved. & G/T\\
29 & Ex.~7.4 & ``metric space generalises Hilbert space'' & More precisely: every \emph{normed} space is a metric space. & T\\
29b & Def.~7.2 & Balls with radius $r\ge0$ & $r>0$ (for $r=0$ the open ball is empty). & T\\
29c & Def.~8.1 & Positive definite kernel: sum $>0$ ``provided that not all $c_i$ are $0$'' & Points $x_1,\dots,x_n$ must be pairwise distinct, otherwise no kernel is positive definite ($x_1=x_2$, $c=(1,-1)$). & M\\
30 & Ex.~8.6, (8.1) & Last equality $\sum_{j,l}c_j\conj{d_l}k(x_j,y_l)=\sum_jc_j\conj{g(x_j)}$ is used without justification. & Requires $k(x,y)=\conj{k(y,x)}$, which holds for every positive semi-definite kernel (Lemma 8.1a, proved). & G\\
31 & Ex.~8.6 & Cauchy--Schwarz used for a form not yet known to be definite; completion ``is again a space of functions'' not proved. & Lemma 1.7a (C--S for semi-definite forms) and full proof of the Moore--Aronszajn theorem (Thm.~8.7) added. & G\\
32 & \S8 & ``we will see a converse \dots\ in the final theorem of this lecture'' -- no theorem is stated. & Stated and proved as Theorem 8.7. & G\\
\bottomrule
\end{longtable}
}
\newpage
%======================================================================
\section{Hilbert spaces}

\begin{definition}\label{def:vs}
Let $X$ be a non-empty set such that one can add its elements (called vectors): there is
a mapping, the vector addition, $X\times X\to X$, $(x,y)\mapsto x+y$, which is
\emph{commutative} ($x+y=y+x$) and \emph{associative} ($(x+y)+z=x+(y+z)$), there is a
neutral element $0$ such that $x+0=x$ for all $x\in X$, and for all $x\in X$ there is an
opposite vector $-x$ such that $x+(-x)=0$. Then $X$ is called a \emph{$\C$-vector space}
if there is another mapping, the scalar multiplication $\C\times X\to X$,
$(\lambda,x)\mapsto\lambda\cdot x$, such that the usual rules of addition and
multiplication are valid: for all $\lambda,\mu\in\C$ and $x,y\in X$,
\[
\lambda\cdot(\mu\cdot x)=(\lambda\mu)\cdot x,\quad 1\cdot x=x,\quad
(\lambda+\mu)\cdot x=\lambda\cdot x+\mu\cdot x,\quad
\lambda\cdot(x+y)=\lambda\cdot x+\lambda\cdot y .
\]
Consequences of these axioms are $0\cdot x=0$ and $(-\lambda)\cdot x=\lambda\cdot(-x)=-(\lambda x)$.

Similarly, $X$ is called an \emph{$\R$-vector space} if there is a scalar multiplication
$\R\times X\to X$ with the same properties.
\end{definition}

\begin{correction}[(Def.~1.1)]
The original definition lists $\lambda(\mu x)=(\lambda\mu)x$, $0\cdot x=0$, $1\cdot x=x$,
$(-\lambda)x=\lambda(-x)$ and $(\lambda+\mu)x=\lambda x+\mu x$, but omits the
commutativity/associativity of $+$ and the distributive law
$\lambda(x+y)=\lambda x+\lambda y$, which are genuine axioms (they do not follow from the
others). On the other hand $0\cdot x=0$ and $(-\lambda)x=-(\lambda x)$ are
\emph{consequences}: $0\cdot x=(0+0)\cdot x=0\cdot x+0\cdot x$, hence $0\cdot x=0$; and
$\lambda x+(-\lambda)x=(\lambda-\lambda)x=0$.
\end{correction}

\begin{example}
Let $n\in\N$. Consider $X=\C^n$ with vector addition
$(x_1,\dots,x_n)+(y_1,\dots,y_n)=(x_1+y_1,\dots,x_n+y_n)$ and scalar multiplication
$\lambda\cdot(x_1,\dots,x_n)=(\lambda x_1,\dots,\lambda x_n)$. Then $X$ is a $\C$-vector
space. Similarly, $X=\R^n$ equipped with formally the same vector addition and scalar
multiplication is an $\R$-vector space.
\end{example}

\begin{example}
Let $a,b\in\R$ such that $a<b$. Let $X=C([a,b])$ be the space of continuous functions
$f:[a,b]\to\C$. If we equip $X$ with vector addition $(f+g)(t)=f(t)+g(t)$ and scalar
multiplication $(\lambda\cdot f)(t)=\lambda f(t)$, then $X$ becomes a $\C$-vector space
(sums and scalar multiples of continuous functions are continuous).
\end{example}

\begin{example}\label{ex:l2}
Let $X=\ell^2=\{(x_n)_{n\in\N}\subseteq\C:\ \sum_{n\in\N}|x_n|^2<\infty\}$. If we equip $X$
with vector addition $(x_n)_n+(y_n)_n=(x_n+y_n)_n$ and scalar multiplication
$\lambda(x_n)_n=(\lambda x_n)_n$, then $X$ becomes a $\C$-vector space.
\end{example}

\begin{proof}
The nontrivial part is to show that the sum of two vectors $(x_n)_n$ and $(y_n)_n$ (which
is clearly a sequence) satisfies $\sum_{n}|x_n+y_n|^2<\infty$. For any $x,y\in\C$ we have
\[
|x+y|^2\leqslant(|x|+|y|)^2\leqslant\big(2\max(|x|,|y|)\big)^2=4\max(|x|^2,|y|^2)
\leqslant4|x|^2+4|y|^2 .
\]
Thus
\[
\sum_{n\in\N}|x_n+y_n|^2\leqslant\sum_{n\in\N}\big(4|x_n|^2+4|y_n|^2\big)
=4\sum_{n\in\N}|x_n|^2+4\sum_{n\in\N}|y_n|^2<\infty .
\]
Closedness under scalar multiplication is clear:
$\sum_n|\lambda x_n|^2=|\lambda|^2\sum_n|x_n|^2<\infty$. All the axioms of
Definition~\ref{def:vs} then hold coordinatewise, since they hold in $\C$.
\end{proof}

\begin{addeddef}{Remark 1.4a (sharper constant)}
Since $0\le(|x|-|y|)^2$ gives $2|x||y|\le|x|^2+|y|^2$, one even has
$|x+y|^2\le|x|^2+2|x||y|+|y|^2\le 2|x|^2+2|y|^2$.
\end{addeddef}

\begin{definition}
Let $X$ be a $\K$-vector space with $\K=\R$ or $\K=\C$. Let $\|\cdot\|:X\to[0,\infty)$ be a
mapping. Then $\|\cdot\|$ is called a \emph{norm} on $X$, and $(X,\|\cdot\|)$ a
\emph{normed vector space}, if
\begin{enumerate}
\item $\|x\|=0\iff x=0$,
\item $\|\lambda x\|=|\lambda|\,\|x\|$ for all $x\in X$ and $\lambda\in\K$ (homogeneity),
\item $\|x+y\|\leqslant\|x\|+\|y\|$ for all $x,y\in X$ (triangle inequality).
\end{enumerate}
\end{definition}

\begin{definition}\label{def:sp}
Let $X$ be a $\K$-vector space with $\K=\R$ or $\K=\C$. A mapping
$\ip{\cdot}{\cdot}:X\times X\to\K$ is called a \emph{scalar product} if
\begin{enumerate}
\item $\ip{x_1+x_2}{y}=\ip{x_1}{y}+\ip{x_2}{y}$ $(x_1,x_2,y\in X)$ and
      $\ip{x}{y_1+y_2}=\ip{x}{y_1}+\ip{x}{y_2}$ $(x,y_1,y_2\in X)$.
\item $\ip{\lambda x}{y}=\conj{\lambda}\ip{x}{y}$ and $\ip{x}{\lambda y}=\lambda\ip{x}{y}$
      $(x,y\in X,\ \lambda\in\K)$.
\item $\ip{x}{y}=\conj{\ip{y}{x}}$ $(x,y\in X)$ (\emph{conjugate symmetry}, also called
      \emph{Hermitian symmetry}).
\item $\ip{x}{x}\geqslant0$, $\ip{x}{x}=0\iff x=0$ $(x\in X)$ (positive definiteness).
\end{enumerate}
Here the properties 1.\ and 2.\ together are called \emph{sesquilinearity} of the scalar product.
\end{definition}

\begin{correction}[(Def.~1.6)]
Property 3 was called ``anti-symetry''. The correct name is \emph{conjugate symmetry}
(or Hermitian symmetry); ``antisymmetry'' usually means $\ip{x}{y}=-\ip{y}{x}$, which is
something different. Note also that 3.\ implies that $\ip{x}{x}\in\R$, so that the
inequality in 4.\ makes sense. For $\K=\R$, the bars are void and 3.\ is ordinary symmetry.
\end{correction}

All our vector spaces below are $\R$- or $\C$-vector spaces. A vector space equipped with a
scalar product is always a normed vector space with a particular norm. This is contained
in the following proposition.

\begin{proposition}\label{prop:CS}
Let $X$ be a vector space equipped with a scalar product $\ip{\cdot}{\cdot}$. Then for all
$x,y\in X$ we have
\begin{enumerate}
\item Cauchy--Schwarz inequality: $|\ip{x}{y}|^2\leqslant\ip{x}{x}\,\ip{y}{y}$.
\item $\|x\|=\sqrt{\ip{x}{x}}$ is a norm on $X$.
\end{enumerate}
In particular, $|\ip{x}{y}|\leqslant\|x\|\,\|y\|$.
\end{proposition}

We prove part 1 in a slightly more general form, which will be needed in Section~8.

\begin{addedthm}{Lemma 1.7a (Cauchy--Schwarz for semi-definite forms)}
Let $X$ be a $\K$-vector space and $s:X\times X\to\K$ satisfy properties 1--3 of
Definition~\ref{def:sp} and $s(x,x)\ge0$ for all $x$ (but not necessarily
$s(x,x)=0\Rightarrow x=0$). Then $|s(x,y)|^2\le s(x,x)\,s(y,y)$ for all $x,y\in X$.
\end{addedthm}
\begin{proof}
Fix $x,y$ and write $c=s(x,y)$. Choose $\omega\in\K$ with $|\omega|=1$ and $c=\omega|c|$.
For $t\in\R$ let $v=y-t\omega x$. By sesquilinearity and conjugate symmetry,
\[
0\le s(v,v)=s(y,y)-t\omega\, s(y,x)-t\conj{\omega}\,s(x,y)+t^2|\omega|^2s(x,x)
= s(y,y)-2t|c|+t^2s(x,x),
\]
because $\omega s(y,x)=\omega\conj{c}=\omega\conj{\omega}|c|=|c|$ and
$\conj\omega s(x,y)=\conj\omega\omega|c|=|c|$.
If $s(x,x)=0$, the affine function $t\mapsto s(y,y)-2t|c|$ is non-negative on $\R$, which
forces $|c|=0$, and the inequality holds. If $s(x,x)>0$, the quadratic polynomial
$t\mapsto s(x,x)t^2-2|c|t+s(y,y)$ is non-negative on $\R$, so its discriminant is
$\le0$: $4|c|^2-4s(x,x)s(y,y)\le0$.
\end{proof}

\begin{proof}[Proof of Proposition~\ref{prop:CS}]
Part 1 is Lemma 1.7a. Part 2 is Exercise~\ref{exo:norm} below. The final inequality is the
square root of part 1.
\end{proof}

\begin{exercise}\label{exo:norm}
Prove the second part of Proposition~\ref{prop:CS}, assuming that the first part is true.
\end{exercise}
\begin{solution}
By property 4 of Definition~\ref{def:sp}, $\ip{x}{x}\ge0$, so $\|x\|=\sqrt{\ip{x}{x}}\in[0,\infty)$ is well defined.
\begin{enumerate}
\item $\|x\|=0\iff\ip{x}{x}=0\iff x=0$ by positive definiteness.
\item $\|\lambda x\|^2=\ip{\lambda x}{\lambda x}=\conj\lambda\lambda\ip{x}{x}=|\lambda|^2\|x\|^2$; take square roots.
\item Using $\ip{y}{x}=\conj{\ip{x}{y}}$, $\Real z\le|z|$ and Cauchy--Schwarz,
\begin{align*}
\|x+y\|^2&=\|x\|^2+\ip{x}{y}+\ip{y}{x}+\|y\|^2=\|x\|^2+2\Real\ip{x}{y}+\|y\|^2\\
&\le\|x\|^2+2\|x\|\|y\|+\|y\|^2=(\|x\|+\|y\|)^2 .\qedhere
\end{align*}
\end{enumerate}
\end{solution}

\begin{definition}
Let $X$ be a normed vector space (e.g.\ a vector space equipped with scalar product and
the norm from Proposition~\ref{prop:CS}). Let $(x_n)_n$ be a sequence in $X$.
\begin{enumerate}
\item The sequence $(x_n)_n$ is called \emph{bounded} if there exists $C>0$ such that
      $\|x_n\|\le C$ for all $n\in\N$.
\item The sequence $(x_n)_n$ is called a \emph{Cauchy sequence} if for all $\epsilon>0$ there
      exists $n_0\in\N$ such that for all $n,m\geqslant n_0$, one has $\|x_n-x_m\|\leqslant\epsilon$.
\item The sequence $(x_n)_n$ is called a \emph{convergent sequence} if there exists some
      $x\in X$ such that for all $\epsilon>0$ there exists $n_0\in\N$ such that for all
      $n\geqslant n_0$, one has $\|x_n-x\|\leqslant\epsilon$.
\end{enumerate}
\end{definition}

\begin{exercise}
Let $(x_n)_n$ be a sequence in a normed vector space.
\begin{enumerate}
\item Show that if $(x_n)_n$ is a Cauchy sequence, then it is bounded.
\item Show that if $(x_n)_n$ is a convergent sequence, then it is Cauchy.
\end{enumerate}
\end{exercise}
\begin{solution}
1. Take $\epsilon=1$: there is $n_0$ with $\|x_n-x_{n_0}\|\le1$ for $n\ge n_0$, hence
$\|x_n\|\le\|x_{n_0}\|+1$ for $n\ge n_0$ (triangle inequality). Then
$C:=\max\{\|x_1\|,\dots,\|x_{n_0}\|\}+1$ satisfies $\|x_n\|\le C$ for all $n$.

2. Let $x$ be the limit and $\epsilon>0$. Choose $n_0$ with $\|x_n-x\|\le\epsilon/2$ for
$n\ge n_0$. Then for $n,m\ge n_0$:
$\|x_n-x_m\|\le\|x_n-x\|+\|x-x_m\|\le\epsilon/2+\epsilon/2=\epsilon$.
\end{solution}

\begin{definition}
Let $X$ be a vector space equipped with a scalar product and the norm from
Proposition~\ref{prop:CS}. If all Cauchy sequences in $X$ converge, then $X$ is called a
\emph{Hilbert space}.
\end{definition}

\begin{example}\label{ex:Cn}
Let $n\in\N$. Let $X=\C^n$ equipped with
\[
\ip{(x_1,\dots,x_n)}{(y_1,\dots,y_n)}=\sum_{k=1}^n\conj{x_k}\,y_k .
\]
Then $\ip{\cdot}{\cdot}$ is a scalar product and $X$ is a Hilbert space. Here, convergence
of a sequence $(x^{(k)})_k=(x_1^{(k)},\dots,x_n^{(k)})_k$ in $X$ simply means that
$\lim_kx_j^{(k)}=x_j$ for all $j=1,\dots,n$ and some limit $x=(x_1,\dots,x_n)$. Similarly,
$X=\R^n$ is a Hilbert space.
\end{example}

\begin{correction}[(Ex.~1.12)]
The original formula reads $\sum_{k=1}^n\conj{x_n}\,y_n$: the summation index $k$ does not
appear in the summand (it would give $n\,\conj{x_n}y_n$). The correct formula is
$\sum_{k=1}^n\conj{x_k}\,y_k$.
\end{correction}

\begin{proof}[Proof of the claims of Example~\ref{ex:Cn}]
Sesquilinearity and conjugate symmetry are immediate from the formula. Moreover
$\ip{x}{x}=\sum_k|x_k|^2\ge0$, with equality iff all $x_k=0$.
\emph{Convergence is coordinatewise:} for every $j$,
$|x_j^{(k)}-x_j|\le\big(\sum_i|x_i^{(k)}-x_i|^2\big)^{1/2}=\|x^{(k)}-x\|$, and conversely
$\|x^{(k)}-x\|^2=\sum_{j=1}^n|x_j^{(k)}-x_j|^2$ is a \emph{finite} sum of terms tending to $0$.
\emph{Completeness:} if $(x^{(k)})_k$ is Cauchy in $\C^n$, then by the same inequality each
coordinate sequence $(x^{(k)}_j)_k$ is Cauchy in $\C$, hence converges to some $x_j$
because $\C$ is complete; by coordinatewise convergence $x^{(k)}\to x=(x_1,\dots,x_n)$.
\end{proof}

\begin{exercise}
Let $X=\C^3$ equipped with the scalar product and induced norm from Example~\ref{ex:Cn}.
Let $x=(2,1,-1)$, $y=(-1,1,i)$, $z=(2,i,0)$ be elements of $X$. Let also
$x_n=\big(e^{-n},\sum_{k=0}^n\frac1{2^k},\frac{n}{n+1}\big)$ for $n\in\N$.
\begin{enumerate}
\item Calculate the scalar products $\ip{x}{y}$, $\ip{y}{z}$ and $\ip{z}{x}$.
\item Calculate the norms $\|x\|,\|y\|,\|z\|$.
\item Determine whether the sequence $(x_n)_n$ converges in $X$ and if it is the case, the limit in $X$.
\end{enumerate}
\end{exercise}
\begin{solution}
Remember that the \emph{first} vector is conjugated.
\begin{enumerate}
\item $\ip{x}{y}=2\cdot(-1)+1\cdot1+(-1)\cdot i=-1-i$;\quad
$\ip{y}{z}=(-1)\cdot2+1\cdot i+\conj{i}\cdot0=-2+i$;\quad
$\ip{z}{x}=2\cdot2+\conj{i}\cdot1+0\cdot(-1)=4-i$.
\item $\|x\|=\sqrt{4+1+1}=\sqrt6$, $\|y\|=\sqrt{1+1+|i|^2}=\sqrt3$, $\|z\|=\sqrt{4+|i|^2+0}=\sqrt5$.
\item Coordinatewise: $e^{-n}\to0$; $\sum_{k=0}^n2^{-k}=\frac{1-2^{-(n+1)}}{1-1/2}=2-2^{-n}\to2$;
$\frac n{n+1}=1-\frac1{n+1}\to1$. By Example~\ref{ex:Cn} the sequence converges in $X$, with
limit $(0,2,1)$. (Explicitly $\|x_n-(0,2,1)\|^2=e^{-2n}+4^{-n}+(n+1)^{-2}\to0$.)\qedhere
\end{enumerate}
\end{solution}

\begin{exercise}
Let $a,b\in\R$, $a<b$. Let $X=C([a,b])$ equipped with
$\ip{f}{g}=\int_a^b\conj{f(t)}\,g(t)\,dt$. Show that this defines a scalar product on $X$.
\end{exercise}
\begin{solution}
The integral exists because $\conj fg$ is continuous on the compact interval $[a,b]$.
Properties 1 and 2 follow from linearity of the integral and from
$\conj{(\lambda f)}=\conj\lambda\,\conj f$. Property 3:
$\conj{\ip{g}{f}}=\conj{\int\conj gf}=\int g\conj f=\ip{f}{g}$
(the integral of the conjugate is the conjugate of the integral). Property 4:
$\ip{f}{f}=\int_a^b|f(t)|^2dt\ge0$. Suppose $\ip{f}{f}=0$ but $f(t_0)\neq0$ for some
$t_0$. By continuity of $|f|^2$ there are $\delta>0$ and an interval $I\subseteq[a,b]$ of
length $\ell>0$ containing $t_0$ with $|f|^2\ge\delta$ on $I$. Then
$\int_a^b|f|^2\ge\delta\ell>0$, a contradiction. Hence $f=0$.
(Here continuity is essential: for Riemann-integrable functions, a function vanishing
except at one point has $\int|f|^2=0$.)
\end{solution}

The space $C([a,b])$ is not a Hilbert space: not every Cauchy sequence converges.

\begin{addition}[(proof that $C({[a,b]})$ is not complete)]
Let $c=\frac{a+b}2$ and, for $n$ large enough that $c+\frac1n\le b$,
\[
f_n(t)=\begin{cases}0,&a\le t\le c,\\ n(t-c),&c\le t\le c+\tfrac1n,\\ 1,&c+\tfrac1n\le t\le b.\end{cases}
\]
Each $f_n$ is continuous. For $m\ge n$, $f_n-f_m$ vanishes outside $[c,c+\frac1n]$ and
$|f_n-f_m|\le1$, so $\|f_n-f_m\|^2\le\frac1n$: $(f_n)$ is a Cauchy sequence.
Suppose $f_n\to f$ for some $f\in C([a,b])$. Since $f_n=0$ on $[a,c]$,
$\int_a^c|f|^2\le\|f-f_n\|^2\to0$, so $f=0$ on $[a,c]$ (by the argument of Exercise~1.14).
For every $\delta>0$ and $n\ge1/\delta$, $f_n=1$ on $[c+\delta,b]$, so
$\int_{c+\delta}^b|f-1|^2\le\|f-f_n\|^2\to0$ and $f=1$ on $[c+\delta,b]$; hence $f=1$ on
$(c,b]$. Then $f$ jumps from $0$ to $1$ at $c$, contradicting continuity. So $(f_n)$ has no
limit in $C([a,b])$. (Its ``limit'' is the indicator function of $(c,b]$; the completion of
$C([a,b])$ is the Lebesgue space $L^2([a,b])$.)
\end{addition}

\begin{example}\label{ex:l2sp}
Let $X=\ell^2$ from Example~\ref{ex:l2}. Equip it with
$\ip{(x_n)_n}{(y_n)_n}=\sum_{n\in\N}\conj{x_n}\cdot y_n$. Then $\ip{\cdot}{\cdot}$ is well
defined and yields a scalar product on $\ell^2$. The space $\ell^2$ is a Hilbert space. It
is the archetype of an infinite dimensional Hilbert space, in the same way as $\C^n$ is the
archetype of a finite ($n$-)dimensional Hilbert space.
\end{example}
\begin{proof}
\emph{Well defined.} By Remark 1.4a,
$|\conj{x_n}y_n|\le\frac12(|x_n|^2+|y_n|^2)$, so the series converges absolutely.
\emph{Scalar product.} Properties 1--3 follow termwise from those in $\C$ (sums of
convergent series are additive, and conjugation is continuous). $\ip{x}{x}=\sum|x_n|^2\ge0$
with equality iff all $x_n=0$.
\emph{Completeness.} Let $(x^{(k)})_k$ be a Cauchy sequence in $\ell^2$, $\epsilon>0$ and
$n_0$ such that $\|x^{(k)}-x^{(l)}\|\le\epsilon$ for $k,l\ge n_0$. For each fixed $n$,
$|x^{(k)}_n-x^{(l)}_n|\le\|x^{(k)}-x^{(l)}\|$, so $(x_n^{(k)})_k$ is Cauchy in $\C$ and has
a limit $x_n$. For each $N\in\N$ and $k,l\ge n_0$,
$\sum_{n=1}^N|x^{(k)}_n-x^{(l)}_n|^2\le\epsilon^2$. Letting $l\to\infty$ (finite sum!)
gives $\sum_{n=1}^N|x^{(k)}_n-x_n|^2\le\epsilon^2$, and then letting $N\to\infty$,
$\sum_{n=1}^\infty|x^{(k)}_n-x_n|^2\le\epsilon^2$ for $k\ge n_0$. Hence
$x^{(k)}-x\in\ell^2$, so $x=x^{(k)}-(x^{(k)}-x)\in\ell^2$, and $\|x^{(k)}-x\|\le\epsilon$
for $k\ge n_0$, i.e.\ $x^{(k)}\to x$ in $\ell^2$.
\end{proof}

\begin{exercise}
Let $X=\ell^2$.
\begin{enumerate}
\item Verify whether the following sequences belong to $X$:
$(x_n)_n=(n)_n$, $(y_n)_n=\big(\frac1{\sqrt{n+1}}\big)_n$, $(z_n)_n=\big(\frac1{2^n}\big)_n$.
\item For each of the above sequences, in case that they belong to $X$, calculate their
norm defined by $\|x\|=\sqrt{\ip{x}{x}}$, where the scalar product is the one from
Example~\ref{ex:l2sp}.
\end{enumerate}
\end{exercise}
\begin{solution}
1. $\sum_n|x_n|^2=\sum_nn^2=\infty$: $(x_n)\notin\ell^2$ (its terms do not even tend to $0$).
$\sum_n|y_n|^2=\sum_n\frac1{n+1}=\infty$ (harmonic series): $(y_n)\notin\ell^2$, although
$y_n\to0$. $\sum_n|z_n|^2=\sum_n4^{-n}<\infty$ (geometric series): $(z_n)\in\ell^2$.

2. With $\N=\{1,2,\dots\}$: $\|z\|^2=\sum_{n=1}^\infty4^{-n}=\frac{1/4}{1-1/4}=\frac13$, so
$\|z\|=\frac1{\sqrt3}$. (If one indexes from $n=0$, then $\|z\|^2=\frac1{1-1/4}=\frac43$ and
$\|z\|=\frac2{\sqrt3}$.)
\end{solution}

\begin{definition}\label{def:bdd}
Let $X,Y$ be Hilbert spaces, equipped with scalar products and the induced norms.
\begin{enumerate}
\item A \emph{bounded linear operator} $T$ is a mapping $T:X\to Y$ that is linear:
$T(x+y)=T(x)+T(y)$, $T(\lambda x)=\lambda T(x)$ for all $x,y\in X$ and $\lambda\in\C$ (or
$\R$), and for which there exists a constant $C>0$ such that $\|T(x)\|\le C\|x\|$ for all
$x\in X$.
\item A \emph{bounded linear form} on $X$ is a bounded linear operator $T:X\to\C$, where
$Y=\C$ is indeed a $\C$-vector space and Hilbert space.
\end{enumerate}
\end{definition}

\begin{remark}\label{rem:finitedim}
In the case that the space $X$ in Definition~\ref{def:bdd} above is finite dimensional, the
boundedness condition is automatic. That is, for any Hilbert space $Y$ and any linear
mapping $T:X\to Y$, there exists $C>0$ such that $\|T(x)\|\le C\|x\|$ for all $x\in X$.
\end{remark}
\begin{proof}
Let $e_1,\dots,e_n$ be an orthonormal basis of $X$ (it exists by Gram--Schmidt, see
Proposition~2.8). Then $x=\sum_j\ip{e_j}{x}e_j$ and $|\ip{e_j}{x}|\le\|x\|$, so
\[
\|T(x)\|=\Big\|\sum_{j=1}^n\ip{e_j}{x}T(e_j)\Big\|\le\sum_{j=1}^n|\ip{e_j}{x}|\,\|T(e_j)\|
\le\Big(\sum_{j=1}^n\|T(e_j)\|\Big)\|x\| .
\]
Take $C=\max\{1,\sum_j\|T(e_j)\|\}$.
\end{proof}

\begin{theorem}[Riesz' Theorem]\label{thm:Riesz}
Let $X$ be a Hilbert space. The bounded linear forms on $X$ are exactly all mappings of
the form $T=T_y:X\to\C$, $T_y(x)=\ip{y}{x}$ with $y\in X$. Moreover $y$ is unique and
$\|T_y\|:=\sup_{\|x\|\le1}|T_y(x)|=\|y\|$.
\end{theorem}

\begin{exercise}\label{exo:Ty}
Show that $T_y$ in Riesz' Theorem~\ref{thm:Riesz} is indeed a bounded linear form. (The
nontrivial part of Theorem~\ref{thm:Riesz} is that every bounded linear form is of this shape.)
\end{exercise}
\begin{solution}
Linearity: $T_y(x_1+\lambda x_2)=\ip{y}{x_1+\lambda x_2}=\ip{y}{x_1}+\lambda\ip{y}{x_2}$
because the scalar product is linear in the second variable. Boundedness: by
Cauchy--Schwarz, $|T_y(x)|\le\|y\|\,\|x\|$, so $C=\|y\|$ works (or $C=1$ if $y=0$).
Moreover $T_y(y/\|y\|)=\|y\|$ for $y\ne0$, hence $\|T_y\|=\|y\|$.
\end{solution}

\begin{addition}[(proof of Riesz' Theorem)]
We need one geometric fact.

\smallskip\noindent\textbf{Lemma 1.19a (orthogonal vector).} \emph{Let $N\subsetneq X$ be a
closed linear subspace of a Hilbert space. Then there is $z\in X$, $z\ne0$, with
$\ip{n}{z}=0$ for all $n\in N$.}

\smallskip\noindent\emph{Proof.} Take $w\in X\setminus N$, let
$d=\inf_{n\in N}\|w-n\|$ and $n_k\in N$ with $\|w-n_k\|^2\to d^2$. The
\emph{parallelogram law} $\|u+v\|^2+\|u-v\|^2=2\|u\|^2+2\|v\|^2$ (expand both sides) with
$u=w-n_k$, $v=w-n_l$ gives, since $\frac{n_k+n_l}2\in N$,
\[
\|n_k-n_l\|^2=2\|w-n_k\|^2+2\|w-n_l\|^2-4\Big\|w-\tfrac{n_k+n_l}2\Big\|^2
\le2\|w-n_k\|^2+2\|w-n_l\|^2-4d^2\xrightarrow[k,l\to\infty]{}0 .
\]
So $(n_k)$ is Cauchy and converges to some $n_0\in X$; $n_0\in N$ as $N$ is closed, and
$\|w-n_0\|=d$. Put $z=w-n_0\ne0$ (as $w\notin N$). For $n\in N$, $t\in\C$, we have
$n_0+tn\in N$, so
$d^2\le\|z-tn\|^2=d^2-2\Real\big(t\ip{z}{n}\big)+|t|^2\|n\|^2$. With
$t=s\,\conj{\ip{z}{n}}$, $s>0$: $0\le-2s|\ip{z}{n}|^2+s^2|\ip{z}{n}|^2\|n\|^2$; divide by
$s$ and let $s\to0$: $\ip{z}{n}=0$. \hfill$\square$

\smallskip\noindent\emph{Proof of Theorem~\ref{thm:Riesz}.} Let $T$ be a bounded linear form.
If $T=0$ take $y=0$. Otherwise $N=\Ker T$ is a linear subspace, $N\ne X$, and $N$ is closed:
if $n_k\in N$, $n_k\to x$, then $|T(x)|=|T(x-n_k)|\le C\|x-n_k\|\to0$. Let $z$ be as in
Lemma~1.19a. For every $x\in X$, the vector $T(x)z-T(z)x$ lies in $N$ (apply $T$), hence
\begin{gather*}
0=\ip{z}{T(x)z-T(z)x}=T(x)\|z\|^2-T(z)\ip{z}{x},\quad\text{hence}\\
T(x)=\frac{T(z)}{\|z\|^2}\ip{z}{x}=\ip{y}{x}\qquad\text{with}\qquad y:=\frac{\conj{T(z)}}{\|z\|^2}\,z .
\end{gather*}
\emph{Uniqueness:} if $\ip{y}{x}=\ip{y'}{x}$ for all $x$, take $x=y-y'$: $\|y-y'\|^2=0$.
The equality $\|T_y\|=\|y\|$ was shown in Exercise~\ref{exo:Ty}. \hfill$\square$

\smallskip In finite dimension the lemma is elementary: complete an orthonormal basis of
$N$ to one of $X$ (Gram--Schmidt) and take one of the new vectors. Then simply
$y=\sum_j\conj{T(e_j)}e_j$ for any orthonormal basis $(e_j)$.
\end{addition}

\begin{remark}\label{rem:braket}
When you will pass to the physics part of this lecture, the notation will typically be the
following: a vector $x$ in a Hilbert space $X$ is denoted $x=\ket{\phi}$ and a bounded
linear form $T_y$ is denoted $T_y=\bra{\psi}$. Then by Riesz' Theorem~\ref{thm:Riesz},
there exists a vector $y\in X$ such that $T_y(x)=\bra{\psi}(\ket{\phi})=\ip{y}{x}$. Thus the
physicist notation is an intuitive way to memorize Riesz' Theorem; the linear form
$\bra{\psi}$ is called \emph{bra}, the vector $\ket{\phi}$ is called \emph{ket}, and when
putting them together as $T_y(x)\in\C$ as above, one obtains a bra-ket, or
\emph{bracket} $\langle\psi|\phi\rangle$.
\end{remark}

\begin{correction}[(notation)]
In the original, bras and kets are typeset with the order relations, $|\phi>$ and
$<\psi|$. The correct symbols are the angle brackets: $\ket{\phi}$, $\bra{\psi}$,
$\langle\psi|\phi\rangle$ (\verb+\langle+, \verb+\rangle+ in \LaTeX).
\end{correction}

\begin{example}\label{ex:ketbra}
Let $X$ be a Hilbert space and $x=\ket{\phi}$ some vector in $X$, as well as $y=\bra{\psi}$
some bounded linear form on $X$. Then $T=\ketbra{\phi}{\psi}$ is a bounded linear operator
$T:X\to X$. It is given, for $z=\ket{\xi}\in X$, by
$T(z)=\ketbra{\phi}{\psi}(\ket{\xi})=\ket{\phi}\langle\psi|\xi\rangle=\ip{y}{z}\,x$.
\end{example}

\begin{exercise}
In the setting of Example~\ref{ex:ketbra} above, show that $T$ is indeed linear $X\to X$
and bounded.
\end{exercise}
\begin{solution}
$T(z_1+\lambda z_2)=\ip{y}{z_1+\lambda z_2}x=\ip{y}{z_1}x+\lambda\ip{y}{z_2}x
=T(z_1)+\lambda T(z_2)$. By Cauchy--Schwarz
$\|T(z)\|=|\ip{y}{z}|\,\|x\|\le\|x\|\,\|y\|\,\|z\|$, so $T$ is bounded with constant
$C=\|x\|\|y\|$ (and in fact $\|T\|=\|x\|\|y\|$, attained at $z=y$).
\end{solution}

\begin{exercise}
Let $X=\C^3$ and $x=\ket{\phi}=(2,-1,0)$, $y=\bra{\psi}=(5,1,-1)$ and $z=\ket{\xi}=(1,0,3)$.
Calculate $\ketbra{\phi}{\psi}(\ket{\xi})$.
\end{exercise}
\begin{solution}
$\langle\psi|\xi\rangle=\ip{y}{z}=5\cdot1+1\cdot0+(-1)\cdot3=2$, hence
$\ketbra{\phi}{\psi}(\ket\xi)=2\,(2,-1,0)=(4,-2,0)$.
\end{solution}
%======================================================================
\section{Tensor products}

\begin{remark}
In the following definition, we recall that a \emph{basis} of a vector space $X$ is some
collection $(x_k)_k\subseteq X$ such that any $x\in X$ can be uniquely written as
$x=\sum_k\lambda_kx_k$ where $\lambda_k\in\C$ and all but finitely many of the $\lambda_k$
are $0$. A fundamental result of linear algebra is that any vector space has a basis. In
case that there is a finite basis $(x_k)_k=x_1,\dots,x_n$, $X$ is called finite
dimensional. In this case any basis consists of $n$ vectors. In this section we are only
concerned with finite dimensional vector spaces.
\end{remark}

\begin{definition}\label{def:tensor}
Let $X$ and $Y$ be finite dimensional vector spaces. Suppose that $X$ has the basis
$x_1,\dots,x_n$ and $Y$ has the basis $y_1,\dots,y_k$. We define the \emph{tensor product}
of $X$ and $Y$, written $X\otimes Y$, as
\[
X\otimes Y=\Big\{\sum_{i=1}^n\sum_{j=1}^k\lambda_{i,j}\,x_i\otimes y_j:\ \lambda_{i,j}\in\C\Big\}.
\]
Here $x_i\otimes y_j$ is first an abstract element. The space $X\otimes Y$ is again a vector
space with addition
\[
\Big(\sum_{i,j}\lambda_{i,j}x_i\otimes y_j\Big)+\Big(\sum_{i,j}\mu_{i,j}x_i\otimes y_j\Big)
=\sum_{i,j}(\lambda_{i,j}+\mu_{i,j})\,x_i\otimes y_j
\]
and scalar multiplication
\[
\lambda\Big(\sum_{i,j}\lambda_{i,j}x_i\otimes y_j\Big)=\sum_{i,j}\lambda\lambda_{i,j}\,x_i\otimes y_j .
\]
The space $X\otimes Y$ possesses naturally the basis $x_i\otimes y_j$, $i=1,\dots,n$,
$j=1,\dots,k$, and thus is again a finite dimensional vector space. Actually
$\dim(X\otimes Y)=\dim(X)\cdot\dim(Y)$.
\end{definition}

\begin{remark}\label{rem:bilin}
If $X$ and $Y$ are finite dimensional vector spaces and $x=\sum_i\lambda_ix_i\in X$ and
$y=\sum_j\mu_jy_j\in Y$, then one writes
$x\otimes y=\sum_{i,j}\lambda_i\mu_j\,x_i\otimes y_j\in X\otimes Y$. One can check that
$x\otimes y+x\otimes z=x\otimes(y+z)$, $x\otimes y+t\otimes y=(x+t)\otimes y$ and
$\lambda(x\otimes y)=(\lambda x)\otimes y=x\otimes(\lambda y)$ for all $x,t\in X$,
$y,z\in Y$ and $\lambda\in\C$. Actually, these observations are typically used to define
the tensor product of not necessarily finite dimensional vector spaces. Any element $p$ of
$X\otimes Y$ can be written as a finite sum $p=\sum_ka_k\otimes b_k$ with the $a_k\in X$ and
the $b_k\in Y$. This way of writing $p$ is not unique in general.
\end{remark}

\begin{addition}[(verifications for Remark~2.3 and independence of the bases)]
\emph{Bilinearity.} With $x=\sum_i\lambda_ix_i$, $y=\sum_j\mu_jy_j$, $z=\sum_j\nu_jy_j$:
$x\otimes y+x\otimes z=\sum_{i,j}\lambda_i(\mu_j+\nu_j)x_i\otimes y_j=x\otimes(y+z)$, since
$y+z$ has coordinates $\mu_j+\nu_j$. The other identities are checked in the same way.
\emph{Finite sums.} $p=\sum_{i,j}\lambda_{i,j}x_i\otimes y_j=\sum_{i,j}(\lambda_{i,j}x_i)\otimes y_j$.
\emph{Non-uniqueness.} E.g.\ $x_1\otimes y_1+x_1\otimes y_2=x_1\otimes(y_1+y_2)$.

\emph{The definition does not depend on the chosen bases.} Let $(u_p)_{p\le n}$ and
$(v_q)_{q\le k}$ be other bases of $X$ and $Y$. Writing $x_i=\sum_p\alpha_{pi}u_p$ and
$y_j=\sum_q\beta_{qj}v_q$, bilinearity gives
$x_i\otimes y_j=\sum_{p,q}\alpha_{pi}\beta_{qj}\,u_p\otimes v_q$. Hence the $nk$ vectors
$u_p\otimes v_q$ span the $nk$-dimensional space $X\otimes Y$, so they form a basis. In
other words, one obtains the same space (with the same elements $x\otimes y$) starting
from any pair of bases.
\end{addition}

\begin{remark}
In the same way as in Definition~\ref{def:tensor}, one can define multiple tensor products:
let $X_1,\dots,X_n$ be finite dimensional vector spaces with bases
$(x^{(1)}_{k_1})_{k_1},\dots,(x^{(n)}_{k_n})_{k_n}$. Then $X_1\otimes X_2\otimes\dots\otimes X_n$
is the vector space
\[
\Big\{\sum_{k_1,\dots,k_n}\lambda_{k_1,\dots,k_n}\,x^{(1)}_{k_1}\otimes\dots\otimes x^{(n)}_{k_n}:\
\lambda_{k_1,\dots,k_n}\in\C\Big\}.
\]
\end{remark}

\begin{remark}
In physicists' notation, the tensor product of $\ket{\psi_1},\dots,\ket{\psi_n}$ belonging
respectively to Hilbert spaces $X_1,\dots,X_n$ is often denoted by
$\ket{\psi_1,\dots,\psi_n}$ (or $\ket{\psi_1}\ket{\psi_2}\cdots\ket{\psi_n}$). For
example, the $n$-qubit space is $(\C^2)^{\otimes n}\cong\C^{2^n}$ with basis
$\ket{b_1b_2\dots b_n}$, $b_i\in\{0,1\}$.
\end{remark}

Particularly relevant for us will be tensor products of (finite dimensional)
\emph{Hilbert} spaces.

\begin{proposition}\label{prop:tensorsp}
Let $X,Y$ be finite dimensional Hilbert spaces. Then $X\otimes Y$ is again a (finite
dimensional) Hilbert space when equipped with the following scalar product. Let
$a=\sum_kx_k\otimes y_k$ and $b=\sum_jx'_j\otimes y'_j$ belong to the tensor product
$X\otimes Y$. Then
\[
\ip{a}{b}_{X\otimes Y}=\sum_{k,j}\ip{x_k}{x'_j}_X\,\ip{y_k}{y'_j}_Y .
\]
Even if the representation of $a$ and $b$ as above is not unique (as we have stated in
Remark~\ref{rem:bilin} above), the expression $\ip{\cdot}{\cdot}_{X\otimes Y}$ above is
well-defined (i.e.\ $\ip{a}{b}_{X\otimes Y}$ is independent of the representations of $a$
and $b$). Moreover, $\ip{\cdot}{\cdot}_{X\otimes Y}$ is a scalar product and $X\otimes Y$
is a Hilbert space when equipped with it.

In the same way, if $X_1,\dots,X_n$ are finite dimensional Hilbert spaces and
$a=\sum_kx_k^{(1)}\otimes\dots\otimes x_k^{(n)}$ and
$b=\sum_ly_l^{(1)}\otimes\dots\otimes y_l^{(n)}$, then
$\ip{a}{b}=\sum_{k,l}\ip{x_k^{(1)}}{y_l^{(1)}}\cdots\ip{x_k^{(n)}}{y_l^{(n)}}$ defines a scalar
product on $X_1\otimes\dots\otimes X_n$.
\end{proposition}
\begin{proof}
(The proof uses orthonormal bases, which exist by Proposition~\ref{prop:onb} below; its
proof does not use Proposition~\ref{prop:tensorsp}.) Let $(e_i)_{i\le n}$ and
$(f_j)_{j\le m}$ be orthonormal bases of $X$ and $Y$. By the Complement above,
$(e_i\otimes f_j)_{i,j}$ is a basis of $X\otimes Y$, so every $a\in X\otimes Y$ has unique
coordinates $\alpha_{ij}$ with $a=\sum_{i,j}\alpha_{ij}e_i\otimes f_j$. Define
\[
[a,b]:=\sum_{i,j}\conj{\alpha_{ij}}\,\beta_{ij}\qquad\text{if }b=\sum_{i,j}\beta_{ij}e_i\otimes f_j .
\]
In coordinates this is the scalar product of $\C^{nm}$ (Example~\ref{ex:Cn}); hence
$[\cdot,\cdot]$ is a scalar product and $X\otimes Y$ is complete for it, i.e.\ a Hilbert
space.

Now let $a=\sum_kx_k\otimes y_k$ be \emph{any} representation. Expanding
$x_k=\sum_i\ip{e_i}{x_k}e_i$ and $y_k=\sum_j\ip{f_j}{y_k}f_j$ (Proposition~\ref{prop:onb})
and using bilinearity, $\alpha_{ij}=\sum_k\ip{e_i}{x_k}\ip{f_j}{y_k}$; similarly
$\beta_{ij}=\sum_l\ip{e_i}{x'_l}\ip{f_j}{y'_l}$. Therefore, by
Proposition~\ref{prop:onb}(5),
\[
[a,b]=\sum_{k,l}\Big(\sum_i\ip{x_k}{e_i}\ip{e_i}{x'_l}\Big)\Big(\sum_j\ip{y_k}{f_j}\ip{f_j}{y'_l}\Big)
=\sum_{k,l}\ip{x_k}{x'_l}_X\,\ip{y_k}{y'_l}_Y .
\]
So the right-hand side of the formula in the proposition equals $[a,b]$ for every choice
of representations: it is well-defined, it is a scalar product, and $X\otimes Y$ is a
Hilbert space. The case of $n$ factors is identical, with orthonormal basis
$e^{(1)}_{i_1}\otimes\dots\otimes e^{(n)}_{i_n}$.
\end{proof}

\begin{exercise}
Let $X=Y=\R^3$. Let $x_1=(1,2,1)$, $y_1=(2,0,-1)$, $x_2=(1,1,2)$, $y_2=(3,0,2)$, and
$a=(1,-1,-1)$, $b=(2,1,0)$.
\begin{enumerate}
\item Justify briefly why $r=x_1\otimes y_1+x_2\otimes y_2$ and $s=a\otimes b$ belong to $X\otimes Y$.
\item Assume now that $X$ and $Y$ are equipped with their standard scalar product and that
$X\otimes Y$ is equipped with the inherited scalar product as in
Proposition~\ref{prop:tensorsp}. Calculate $\ip{r}{s}_{X\otimes Y}$.
\end{enumerate}
\end{exercise}
\begin{solution}
1. By Remark~\ref{rem:bilin}, $x\otimes y\in X\otimes Y$ for all $x\in X$, $y\in Y$
(expand $x$, $y$ in the standard bases); $X\otimes Y$ is a vector space, so the sum
$r$ also belongs to it.

2. By Proposition~\ref{prop:tensorsp},
$\ip{r}{s}=\ip{x_1}{a}\ip{y_1}{b}+\ip{x_2}{a}\ip{y_2}{b}$ with
$\ip{x_1}{a}=1-2-1=-2$, $\ip{y_1}{b}=4+0+0=4$, $\ip{x_2}{a}=1-1-2=-2$, $\ip{y_2}{b}=6+0+0=6$.
Hence $\ip{r}{s}_{X\otimes Y}=(-2)\cdot4+(-2)\cdot6=-20$.
(Check with Kronecker vectors in $\R^9$, see Exercise 2.11:
$r=(5,0,1,7,0,0,8,0,3)$, $s=(2,1,0,-2,-1,0,-2,-1,0)$ and $r\cdot s=10-14-16=-20$.)
\end{solution}

Hilbert spaces do not only possess a basis but also an \emph{orthonormal} basis. What this
is, see the next proposition.

\begin{proposition}\label{prop:onb}
Let $X$ be a Hilbert space. Then there exists a family $(x_k)_k$ of linearly independent
vectors in $X$ that have the following properties:
\begin{enumerate}
\item The vectors $x_k$ are normalized: $\|x_k\|=1$ for all $k$.
\item The vectors $x_k$ are orthogonal to each other: $\ip{x_k}{x_l}=0$ for $k\ne l$.
\item Any $x\in X$ can be written as a convergent series $x=\sum_k\lambda_kx_k$ with unique
coefficients $\lambda_k=\ip{x_k}{x}$.
\item The norm of any $x\in X$ can be computed as $\|x\|^2=\sum_k|\lambda_k|^2=\sum_k|\ip{x_k}{x}|^2$.
\item The scalar product of two vectors $x,y\in X$ can be computed as
$\ip{x}{y}=\sum_k\ip{x}{x_k}\ip{x_k}{y}$.
\end{enumerate}
A family with the above properties is called \emph{orthonormal basis}.
\end{proposition}
\begin{proof}[Proof (finite dimensional case)]
Let $v_1,\dots,v_n$ be a basis of $X$. \emph{Gram--Schmidt:} put $u_1=v_1/\|v_1\|$ and,
recursively for $k=2,\dots,n$,
\[
w_k=v_k-\sum_{j<k}\ip{u_j}{v_k}u_j,\qquad u_k=w_k/\|w_k\| .
\]
$w_k\neq0$, since otherwise $v_k\in\spn\{u_1,\dots,u_{k-1}\}=\spn\{v_1,\dots,v_{k-1}\}$,
contradicting linear independence. By induction, $u_k$ is normalized and for $l<k$:
\[
\ip{u_l}{w_k}=\ip{u_l}{v_k}-\sum_{j<k}\ip{u_j}{v_k}\ip{u_l}{u_j}=\ip{u_l}{v_k}-\ip{u_l}{v_k}=0 .
\]
Orthonormal vectors are linearly independent (if $\sum_k\lambda_ku_k=0$, take the scalar
product with $u_l$: $\lambda_l=0$); $n$ of them form a basis. This gives 1 and 2.
3.: write $x=\sum_k\lambda_ku_k$ and take the scalar product with $u_l$:
$\ip{u_l}{x}=\sum_k\lambda_k\ip{u_l}{u_k}=\lambda_l$.
5.: with $\mu_l=\ip{u_l}{y}$,
\[
\ip{x}{y}=\Big\langle\sum_k\lambda_ku_k,\sum_l\mu_lu_l\Big\rangle
=\sum_{k,l}\conj{\lambda_k}\mu_l\delta_{kl}=\sum_k\conj{\ip{u_k}{x}}\ip{u_k}{y}
=\sum_k\ip{x}{u_k}\ip{u_k}{y}.
\]
4.\ is 5.\ with $y=x$ (Parseval's identity).
\end{proof}

\begin{addition}[(infinite dimensional case)]
If $X$ is separable (has a countable dense subset $\{v_1,v_2,\dots\}$), apply Gram--Schmidt
to the $v_k$, discarding those already in the span of the previous ones; the resulting
orthonormal family $(u_k)$ has dense span. \emph{Bessel's inequality}
$\sum_k|\ip{u_k}{x}|^2\le\|x\|^2$ (which follows from
$0\le\|x-\sum_{k\le N}\ip{u_k}{x}u_k\|^2=\|x\|^2-\sum_{k\le N}|\ip{u_k}{x}|^2$) shows that
$\sum_k\ip{u_k}{x}u_k$ is Cauchy, hence converges in $X$ (completeness!) to some $x'$ with
$x-x'\perp u_k$ for all $k$, hence $x-x'\perp\overline{\spn}\{u_k\}=X$, i.e.\ $x=x'$. For
non-separable $X$ one uses Zorn's lemma to get a maximal orthonormal family, which is then
uncountable; the ``series'' in 3.--5.\ must be understood as unconditional sums, in which
for each $x$ only countably many coefficients $\ip{x_k}{x}$ are non-zero (by Bessel).
\end{addition}

\begin{remark}
In case that $X$ is an infinite dimensional Hilbert space, its orthonormal bases are
infinite as well. In that case, an orthonormal basis is not a basis in the sense of linear
algebra since in Proposition~\ref{prop:onb} above, the sum in the third property has in
general infinitely many non-zero coefficients $\lambda_k$, as opposed to a basis, where the
sum has always only finitely many non-zero $\lambda_k$. But for the rest of this lecture,
this ambiguity is not important, since we are mainly concerned with finite dimensional
Hilbert spaces, in which case indeed an orthonormal basis is always a basis.
\end{remark}

How do orthonormal bases behave when we take tensor products? The answer is quite
satisfactory as reveals the following proposition.

\begin{proposition}\label{prop:tensoronb}
Let $X,Y$ be finite dimensional Hilbert spaces. Let $(x_k)_k$ and $(y_l)_l$ be orthonormal
bases of $X$ and $Y$ respectively. Then $(x_k\otimes y_l)_{k,l}$ is an orthonormal basis of
$X\otimes Y$. In particular, for $a,b\in X\otimes Y$, the scalar product
$\ip{a}{b}_{X\otimes Y}$ can be calculated as
\[
\ip{a}{b}_{X\otimes Y}=\sum_{k,l}\ip{a}{x_k\otimes y_l}\ip{x_k\otimes y_l}{b} .
\]
Moreover, if $a=\sum_{k,l}\lambda_{k,l}\,x_k\otimes y_l$ belongs to $X\otimes Y$, then
$\|a\|^2=\sum_{k,l}|\lambda_{k,l}|^2$.
\end{proposition}
\begin{proof}
By Proposition~\ref{prop:tensorsp},
$\ip{x_k\otimes y_l}{x_{k'}\otimes y_{l'}}=\ip{x_k}{x_{k'}}\ip{y_l}{y_{l'}}=\delta_{kk'}\delta_{ll'}$,
so the family is orthonormal. It has $\dim X\cdot\dim Y=\dim(X\otimes Y)$ elements and is
linearly independent, hence it is a basis, and thus an orthonormal basis. The two formulas
are properties 5 and 4 of Proposition~\ref{prop:onb}, with
$\lambda_{k,l}=\ip{x_k\otimes y_l}{a}$.
\end{proof}

\begin{exercise}
Let $X=\R^2$ and $Y=\R^3$ equipped with the orthonormal bases $x_1=(1,0)$, $x_2=(0,1)$, as
well as $y_1=(1,0,0)$, $y_2=(0,1,0)$, $y_3=(0,0,1)$. Find a convenient way to write the
resulting orthonormal basis of $X\otimes Y=\R^2\otimes\R^3$.
\end{exercise}
\begin{solution}
By Proposition~\ref{prop:tensoronb}, $(x_i\otimes y_j)_{i\le2,j\le3}$ is an orthonormal
basis of the 6-dimensional space $\R^2\otimes\R^3$. Two convenient identifications:

\emph{(a) Kronecker vectors.} Identify $\R^2\otimes\R^3\cong\R^6$ via
$a\otimes b\mapsto(a_1b_1,a_1b_2,a_1b_3,a_2b_1,a_2b_2,a_2b_3)=(a_1b,\,a_2b)$. Then
$x_i\otimes y_j\mapsto e_{3(i-1)+j}$:
\[
x_1\otimes y_1=e_1,\ x_1\otimes y_2=e_2,\ x_1\otimes y_3=e_3,\
x_2\otimes y_1=e_4,\ x_2\otimes y_2=e_5,\ x_2\otimes y_3=e_6,
\]
the standard basis of $\R^6$. This identification preserves scalar products since
$\sum_{i,j}a_ib_ja'_ib'_j=\ip{a}{a'}\ip{b}{b'}$. This is the convention used for qubits:
$\ket{0}\otimes\ket{1}=\ket{01}$, etc.

\emph{(b) Matrices.} Identify $a\otimes b$ with the $2\times3$ matrix $ab^T=(a_ib_j)_{ij}$;
then $x_i\otimes y_j$ is the matrix unit $E_{ij}$ (a $1$ at position $(i,j)$, $0$
elsewhere), and the scalar product becomes the Hilbert--Schmidt product $\tr(A^TB)$.
\end{solution}
%======================================================================
\section{Probabilities}

\begin{definition}\label{def:prob}
Let $\Omega$ be some non-empty set, either finite or countably infinite. Let $\PP$ be a
mapping defined on the subsets of $\Omega$ and with values in $[0,1]$. Then $(\Omega,\PP)$
is called a \emph{probability space} if
\begin{enumerate}
\item $\PP(\emptyset)=0$, $\PP(\Omega)=1$,
\item $\PP\big(\bigcup_{n\in\N}A_n\big)=\sum_{n\in\N}\PP(A_n)$ for any countable family of
pairwise disjoint $A_n$.
\end{enumerate}
The set $\Omega$ is called \emph{universe}, the subsets $A\subseteq\Omega$ are called
\emph{events} and $\PP(A)$ is called its \emph{probability}. Naturally,
$\PP(A)\in[0,1]=[0\%,100\%]$.
\end{definition}

\begin{correction}[(Def.~3.1)]
The original allows $\Omega$ ``finite or infinite''. For \emph{uncountable} $\Omega$ (e.g.\
$\Omega=[0,1]$) it is in general impossible to define $\PP$ on \emph{all} subsets (e.g.\
there is no translation-invariant $\sigma$-additive extension of length to all subsets of
$[0,1]$, by Vitali's construction); one has to restrict to a $\sigma$-algebra of events.
Since Definition~3.11 assumes anyway that $\Omega$ is at most countable ``as is always the
case in this section'', we state this assumption from the start.
\end{correction}

\begin{addition}[(elementary consequences)]
For a countable $\Omega$, put $p_\omega=\PP(\{\omega\})$. Since every $A\subseteq\Omega$ is
the countable disjoint union of its singletons, $\PP(A)=\sum_{\omega\in A}p_\omega$, and
$\sum_{\omega\in\Omega}p_\omega=1$. Conversely, any family $p_\omega\ge0$ with
$\sum_\omega p_\omega=1$ defines a probability by this formula ($\sigma$-additivity is the
rearrangement theorem for series with non-negative terms). Further:
$\PP(\Omega\setminus A)=1-\PP(A)$, and $A\subseteq B\Rightarrow\PP(A)\le\PP(B)$.
\end{addition}

\begin{definition}
Let $(\Omega,\PP)$ be a probability space and $X:\Omega\to\R$ or $X:\Omega\to\C$ be a
function. Then $X$ is called a \emph{random variable} (either real valued (mostly) or
complex valued).
\end{definition}

\begin{example}
Let $\Omega=\{1,\dots,6\}$ and $\PP(A)=\frac16\#A$, where $\#A$ denotes the cardinality of
the subset $A$ of $\Omega$. Then $(\Omega,\PP)$ is a probability space. The function
$X:\Omega\to\R$, $x\mapsto x$ is a random variable designed to implement a rolling dice.
\end{example}

\begin{example}\label{ex:binom}
Let $(\Omega,\PP)$ be a probability space. Let $n\in\N$ and $p\in[0,1]$. Let
$X:\Omega\to\{0,1,\dots,n\}$ be a random variable. Then $X$ is called a \emph{binomial}
variable with parameters $n$ and $p$ if for $k=0,\dots,n$, the sets
$A_k=\{\omega\in\Omega:\ X(\omega)=k\}$ have the probability
\[
\PP(A_k)=\binom nkp^k(1-p)^{n-k}.
\]
For $n=1$ (values in $\{0,1\}$, $\PP(X=1)=p$) $X$ is called a \emph{Bernoulli} variable.
\end{example}

\begin{correction}[(Ex.~3.4)]
The original calls this a ``Bernoulli variable with parameters $n$ and $p$''. This is the
\emph{binomial distribution} $\mathcal B(n,p)$: it counts the number of successes in $n$
independent Bernoulli trials with success probability $p$. A Bernoulli variable is the
special case $n=1$. (The convention $0^0=1$ is used for $p\in\{0,1\}$.)
\end{correction}

\begin{example}
Let $(\Omega,\PP)$ be a probability space. Let $\lambda>0$. Let $X:\Omega\to\N_0$ be a
random variable. Then $X$ is called a \emph{Poisson variable} with parameter $\lambda$ if
for any $k\geqslant0$, the sets $A_k=\{\omega\in\Omega:\ X(\omega)=k\}$ have the probability
\[
\PP(A_k)=e^{-\lambda}\frac{\lambda^k}{k!}.
\]
\end{example}

\begin{exercise}
Show that for a binomial variable $X$ with parameters $n$ and $p$ (``Bernoulli variable'' in
the original), one has $\PP(\{\omega\in\Omega:\ X(\omega)\in\{0,\dots,n\}\})=1$.
Show that for a Poisson variable $X$ with parameter $\lambda$, one has
$\PP(\{\omega\in\Omega:\ X(\omega)\geqslant0\})=1$.
\end{exercise}

\begin{correction}[(Ex.~3.6)]
As literally stated, the exercise is trivial: since $X$ takes its values in $\{0,\dots,n\}$
(resp.\ $\N_0$), the event in question is all of $\Omega$, whose probability is $1$ by
definition. The intended (and instructive) statement is that the probabilities of the
events $A_k$ add up to $1$, i.e.\ that the definitions in Examples 3.4 and 3.5 are
\emph{consistent}: no variable with these properties could exist otherwise.
\end{correction}
\begin{solution}
The sets $A_k$ are pairwise disjoint and $\bigcup_kA_k=\{X\in\{0,\dots,n\}\}=\Omega$.
By $\sigma$-additivity and the binomial theorem,
\[
\PP(\Omega)=\sum_{k=0}^n\PP(A_k)=\sum_{k=0}^n\binom nkp^k(1-p)^{n-k}=(p+(1-p))^n=1 .
\]
For the Poisson variable, $\{X\ge0\}=\bigcup_{k\ge0}A_k$ (disjoint), and by the exponential series,
\[
\sum_{k=0}^\infty\PP(A_k)=e^{-\lambda}\sum_{k=0}^\infty\frac{\lambda^k}{k!}=e^{-\lambda}e^{\lambda}=1 .\qedhere
\]
\end{solution}

\begin{remark}
Note that an easy way to implement a binomial or Poisson variable is as follows. For the
binomial variable, take $\Omega=\{0,\dots,n\}$ with
$\PP(\{k\})=\binom nkp^k(1-p)^{n-k}$ for $k=0,\dots,n$ and $\PP(A)=\sum_{k\in A}\PP(\{k\})$.
Then define $X(\omega)=\omega$. For the Poisson variable, take $\Omega=\N_0$ with
$\PP(\{k\})=e^{-\lambda}\frac{\lambda^k}{k!}$ for $k\in\Omega$. Then define again
$X(\omega)=\omega$. (These are probability spaces by Exercise~3.6 and the Complement after
Definition~\ref{def:prob}.)
\end{remark}

\begin{definition}
Let $(\Omega,\PP)$ be a probability space. Let $A,B\subseteq\Omega$. Then $A$ and $B$ are
called (stochastically) \emph{independent} if $\PP(A\cap B)=\PP(A)\cdot\PP(B)$. ``The
knowledge if $A$ occurs does not lower or raise the probability that $B$ occurs.'' Moreover,
let $X,Y$ be random variables over $\Omega$. Then $X,Y$ are called independent if the sets
$A=\{\omega\in\Omega:\ X(\omega)\in A_X\}$ and $B=\{\omega\in\Omega:\ Y(\omega)\in B_Y\}$ are
independent for all $A_X,B_Y\subseteq\R$ (or $\subseteq\C$ if $X,Y$ are complex valued).
\end{definition}

\begin{exercise}
Let $\Omega=\{1,2,3,4,5,6\}$ together with $\PP(A)=\frac16\#A$ be the universe of rolling a dice.
\begin{enumerate}
\item Let $A=\{2,3,5\}$ be the event ``the dice shows a prime number'', and $B=\{3,6\}$ be
the event ``the dice shows a number divisible by 3''. Determine whether $A$ and $B$ are independent.
\item Let $A$ be the same event as before and $B=\{1,2,3\}$ be the event ``the dice shows a
number smaller or equal than 3''. Determine whether $A$ and $B$ are independent.
\end{enumerate}
\end{exercise}
\begin{solution}
1. $\PP(A)=\frac36=\frac12$, $\PP(B)=\frac26=\frac13$, $A\cap B=\{3\}$,
$\PP(A\cap B)=\frac16=\frac12\cdot\frac13$: \emph{independent}.

2. $\PP(B)=\frac12$, $A\cap B=\{2,3\}$, $\PP(A\cap B)=\frac13\ne\frac14=\PP(A)\PP(B)$:
\emph{not independent} (knowing that the result is $\le3$ raises the probability of a
prime from $\frac12$ to $\frac23$).
\end{solution}

\begin{remark}\label{rem:product}
Let $(\Omega,\PP)$ be a probability space and $X,Y$ be random variables. A standard way to
obtain random variables $X',Y'$ which have the same behaviour as $X,Y$ but which are
independent is as follows. Define $\Omega_2=\Omega\times\Omega$ and
$\PP_2(A\times B)=\PP(A)\PP(B)$ for any $A\times B\subseteq\Omega_2$. Extend $\PP_2$ to a
probability by taking
$\PP_2\big(\bigcup_{n\in\N}A_n\times B_n\big)=\sum_{n\in\N}\PP_2(A_n\times B_n)$ whenever
the $A_n\times B_n$ are pairwise disjoint. Then let $X'(\omega_1,\omega_2)=X(\omega_1)$
and $Y'(\omega_1,\omega_2)=Y(\omega_2)$.
\end{remark}

\begin{addition}[(details for Remark~3.10)]
Since $\Omega$ is countable, so is $\Omega_2$, and every $C\subseteq\Omega_2$ is the
countable disjoint union of the ``rectangles'' $\{\omega_1\}\times\{\omega_2\}$,
$(\omega_1,\omega_2)\in C$. Hence the extension is necessarily
$\PP_2(C)=\sum_{(\omega_1,\omega_2)\in C}p_{\omega_1}p_{\omega_2}$, which is a probability
(Complement after Def.~\ref{def:prob}, since $\sum_{\omega_1,\omega_2}p_{\omega_1}p_{\omega_2}=1$),
and indeed $\PP_2(A\times B)=\sum_{\omega_1\in A}p_{\omega_1}\sum_{\omega_2\in B}p_{\omega_2}=\PP(A)\PP(B)$.

\emph{Same behaviour:} $\{X'\in A_X\}=\{X\in A_X\}\times\Omega$, so
$\PP_2(X'\in A_X)=\PP(X\in A_X)\PP(\Omega)=\PP(X\in A_X)$; similarly for $Y'$.

\emph{Independence:} $\{X'\in A_X\}\cap\{Y'\in B_Y\}=\{X\in A_X\}\times\{Y\in B_Y\}$, so
$\PP_2(X'\in A_X,\,Y'\in B_Y)=\PP(X\in A_X)\PP(Y\in B_Y)=\PP_2(X'\in A_X)\PP_2(Y'\in B_Y)$.
\end{addition}

\begin{definition}\label{def:expect}
Let $X$ be a random variable over the probability space $(\Omega,\PP)$. Assume that
$\Omega$ is at most countable as is always the case in this section. Then
\[
\E(X)=\sum_{\omega\in\Omega}X(\omega)\PP(\{\omega\})
\]
is called the \emph{expectation} of $X$, whenever the series converges \emph{absolutely},
i.e.\ if
\[
\E(|X|)=\sum_{\omega\in\Omega}|X(\omega)|\,\PP(\{\omega\})<\infty
\]
(e.g.\ if $X$ is bounded).
The expectation has the following properties:
\begin{enumerate}
\item $\E(1)=1$,
\item $\E(X+Y)=\E(X)+\E(Y)$ whenever $\E(X)$ and $\E(Y)$ are meaningful,
\item $\E(\lambda X)=\lambda\E(X)$ whenever $\E(X)$ is meaningful,
\item $\E(X)\geqslant0$ whenever $X\geqslant0$,
\item $\E(XY)=\E(X)\E(Y)$ whenever $X,Y$ are independent and their expectations are meaningful.
\end{enumerate}
\end{definition}

\begin{correction}[(Def.~3.11)]
The original requires only that ``the series converges''. But a sum indexed by a countable
set $\Omega$ has no preferred order, and a merely conditionally convergent series can be
rearranged to give any value (Riemann's rearrangement theorem). The correct condition is
\emph{absolute} convergence, under which the value does not depend on the enumeration of
$\Omega$, and which is needed for properties 2 and 5.
\end{correction}

\begin{proof}[Proof of properties 1--5]
1.\ $\E(1)=\sum_\omega p_\omega=1$. 2.\ and 3.: $|X+Y|\le|X|+|Y|$ and $|\lambda X|=|\lambda||X|$
give absolute convergence, and sums/scalar multiples of absolutely convergent series are
computed termwise. 4.: all terms are $\ge0$.

5.\ \emph{Grouping by values.} $X(\Omega)$ and $Y(\Omega)$ are countable. Because of
absolute convergence, we may regroup the terms of the series:
\begin{align*}
\E(XY)&=\sum_{\omega}X(\omega)Y(\omega)p_\omega
=\sum_{a\in X(\Omega)}\sum_{b\in Y(\Omega)}ab\,\PP(X=a,\,Y=b)\\
&=\sum_{a}\sum_{b}ab\,\PP(X=a)\PP(Y=b)=\E(X)\E(Y),
\end{align*}
using independence with $A_X=\{a\}$, $B_Y=\{b\}$, and $\E(X)=\sum_aa\PP(X=a)$ (again by
grouping). The same computation with $|a||b|$ shows $\E(|XY|)=\E(|X|)\E(|Y|)<\infty$, which
justifies the regrouping.
\end{proof}

\begin{exercise}
\begin{enumerate}
\item Let $\Omega=\{1,2,3,4,5,6\}$ with probabilities $\PP(A)=\frac16\#A$ for any
$A\subseteq\Omega$. Let $X:\Omega\to\R$, $X(\omega)=\omega$. Calculate the expectation of $X$.
\item Let $X$ be a binomial random variable with parameters $n\in\N$ and $p\in[0,1]$.
Calculate its expectation.
\item Let $X$ be a Poisson random variable with parameter $\lambda>0$. Calculate its expectation.
\end{enumerate}
\end{exercise}
\begin{solution}
1. $\E(X)=\frac16(1+2+3+4+5+6)=\frac{21}6=\frac72$.

2. By grouping, $\E(X)=\sum_{k=0}^nk\binom nkp^k(1-p)^{n-k}$. Using
$k\binom nk=n\binom{n-1}{k-1}$ for $k\ge1$ and the binomial theorem,
\[
\E(X)=np\sum_{k=1}^n\binom{n-1}{k-1}p^{k-1}(1-p)^{(n-1)-(k-1)}=np\,(p+1-p)^{n-1}=np .
\]
(Alternatively: $X$ has the same distribution as a sum of $n$ independent Bernoulli($p$)
variables, each with expectation $p$; use property 2.)

3. $\E(X)=\sum_{k=0}^\infty ke^{-\lambda}\frac{\lambda^k}{k!}
=\lambda e^{-\lambda}\sum_{k=1}^\infty\frac{\lambda^{k-1}}{(k-1)!}=\lambda e^{-\lambda}e^\lambda=\lambda$.
(The series has non-negative terms, so absolute convergence is automatic. For $\lambda=0$,
which the original also allows, $X=0$ almost surely and $\E(X)=0=\lambda$.)
\end{solution}
%======================================================================
\section{Operators acting on Hilbert space}

Let $X,Y$ be Hilbert spaces. We recall that we have defined the notion of a bounded linear
operator $T:X\to Y$ in Definition~\ref{def:bdd}. As a special case, we examine now what can
be said if $X=\C^n$ and $Y=\C^m$.

\begin{proposition}\label{prop:matrix}
Let $n,m\in\N$ and $X=\C^n$ and $Y=\C^m$ be the Hilbert spaces equipped with the scalar
product from Section~1.
\begin{enumerate}
\item Let $A=(a_{ij})$ be a matrix with $m$ rows and $n$ columns. Then $A$ defines a
(bounded) linear operator $T_A:\C^n\to\C^m$ by the formula
$T_A(x_1,\dots,x_n)=\big(\sum_{j=1}^na_{ij}x_j\big)_{i=1}^m$.
\item Conversely each (bounded) linear operator $T:\C^n\to\C^m$ gives rise to a matrix
$A=(a_{ij})$ with $m$ rows and $n$ columns with
\[
a_{ij}=\ip{f_i}{T(e_j)} ,
\]
where $e_j$ (resp.\ $f_i$) is the vector from the standard basis of $\C^n$ (resp.\ $\C^m$)
having zero entries everywhere except at the $j$-th (resp.\ $i$-th) position where there
is a $1$. Then $T=T_A$ with $T_A$ as in the first case.
\item The correspondence $A\longleftrightarrow T_A$ between matrices and linear operators is
bijective.
\end{enumerate}
\end{proposition}

\begin{correction}[(Prop.~4.1(2))]
The original formula is $a_{ij}=\ip{T(e_j)}{f_i}$. With our scalar product (anti-linear in
the first variable), $\ip{T(e_j)}{f_i}=\conj{(T(e_j))_i}=\conj{a_{ij}}$, so the original
formula gives the complex \emph{conjugate} of the correct entry (e.g.\ for $T=i\cdot\mathrm{Id}$
on $\C^1$ it gives $-i$ instead of $i$). The correct formula is $a_{ij}=\ip{f_i}{T(e_j)}$,
i.e.\ the $i$-th coordinate of the $j$-th column $T(e_j)$. (In mathematics texts where the
scalar product is linear in the first variable, the original formula is correct; the
point is to be consistent with Definition~1.6.) Also: ``lines'' $\to$ ``rows'', and
``bijective (i.e.\ exact)'' $\to$ ``bijective''.
\end{correction}

\begin{proof}
1. Linearity of $T_A$ is clear. Boundedness: by Cauchy--Schwarz in $\C^n$,
\[
\|T_Ax\|^2=\sum_i\Big|\sum_ja_{ij}x_j\Big|^2\le\sum_i\Big(\sum_j|a_{ij}|^2\Big)\|x\|^2,
\]
so $\|T_Ax\|\le\|A\|_{HS}\|x\|$ with $\|A\|_{HS}=(\sum_{i,j}|a_{ij}|^2)^{1/2}$ (or use
Remark~\ref{rem:finitedim}).

2. Write $T(e_j)=\sum_i\ip{f_i}{T(e_j)}f_i=\sum_ia_{ij}f_i$ (Proposition~\ref{prop:onb}).
For $x=\sum_jx_je_j$, linearity gives
$T(x)=\sum_jx_jT(e_j)=\sum_i\big(\sum_ja_{ij}x_j\big)f_i=T_A(x)$.

3. Part 2 shows that $A\mapsto T_A$ is surjective. It is injective: if $T_A=T_B$ then
$a_{ij}=\ip{f_i}{T_A(e_j)}=\ip{f_i}{T_B(e_j)}=b_{ij}$.
\end{proof}

\begin{exercise}
Let $X=Y=\C^3$. Let $\ket{\phi}=(1,2,-1)$ and $\bra{\psi}=(0,3,2)$. Calculate the matrix of
the linear operator $T=\ketbra{\phi}{\psi}$ from Example~\ref{ex:ketbra}.
\end{exercise}
\begin{solution}
$T(z)=\ip{\psi}{z}\phi$, so $T(e_j)=\conj{\psi_j}\,\phi$ and
$a_{ij}=\ip{e_i}{T(e_j)}=\phi_i\conj{\psi_j}$: the matrix of $\ketbra{\phi}{\psi}$ is
$\phi\,\psi^*$ (column times conjugated row). Here $\psi$ is real, so
\[
T=\begin{pmatrix}1\\2\\-1\end{pmatrix}\begin{pmatrix}0&3&2\end{pmatrix}
=\begin{pmatrix}0&3&2\\0&6&4\\0&-3&-2\end{pmatrix}.\qedhere
\]
\end{solution}

\begin{definition}\label{def:adjoint}
Let $X,Y$ be Hilbert spaces and $T:X\to Y$ a bounded linear operator. Then there exists
exactly one bounded linear operator $S:Y\to X$ such that $\ip{Tx}{y}_Y=\ip{x}{Sy}_X$ for all
$x\in X$ and $y\in Y$. This operator $S$ is denoted by $T^*$ and called the \emph{adjoint} of $T$.
\end{definition}

\begin{addition}[(existence and uniqueness of the adjoint)]
Fix $y\in Y$. The map $x\mapsto\ip{y}{Tx}_Y$ is linear and $|\ip{y}{Tx}|\le\|y\|\|T\|\|x\|$,
so it is a bounded linear form on $X$. By Riesz' Theorem~\ref{thm:Riesz} there is a unique
$z=:Sy\in X$ with $\ip{y}{Tx}=\ip{Sy}{x}$ for all $x$; conjugating,
$\ip{Tx}{y}=\ip{x}{Sy}$. \emph{Linearity of $S$:}
$\ip{x}{S(y_1+\lambda y_2)}=\ip{Tx}{y_1+\lambda y_2}=\ip{x}{Sy_1+\lambda Sy_2}$ for all $x$,
so $S(y_1+\lambda y_2)=Sy_1+\lambda Sy_2$ (uniqueness in Riesz).
\emph{Boundedness:} $\|Sy\|^2=\ip{Sy}{Sy}=\ip{TSy}{y}\le\|T\|\|Sy\|\|y\|$, hence
$\|Sy\|\le\|T\|\|y\|$. \emph{Uniqueness:} if also $\ip{Tx}{y}=\ip{x}{S'y}$, then
$\ip{x}{(S-S')y}=0$ for all $x$; take $x=(S-S')y$.
\end{addition}

\begin{exercise}\label{exo:adjmatrix}
Let $A=(a_{ij})$ be an $m\times n$ matrix and $T=T_A:\C^n\to\C^m$ be the associated linear
operator. Determine the matrix of $T^*$ in the sense of Proposition~\ref{prop:matrix}.
\end{exercise}
\begin{solution}
$T^*:\C^m\to\C^n$ has an $n\times m$ matrix $B=(b_{ij})$ with, by Proposition~\ref{prop:matrix},
\[
b_{ij}=\ip{e_i}{T^*f_j}=\ip{Te_i}{f_j}=\conj{\ip{f_j}{Te_i}}=\conj{a_{ji}} .
\]
So the matrix of $T^*$ is the \emph{conjugate transpose} $A^*=\conj{A}^{\,T}=(\conj{a_{ji}})_{ij}$.
(Note that with the uncorrected formula of Prop.~4.1 one would have to be very careful
here: the correct formula is essential.) Check: $\ip{Ax}{y}=\sum_{i,j}\conj{a_{ij}x_j}\,y_i
=\sum_j\conj{x_j}\sum_i\conj{a_{ij}}y_i=\ip{x}{A^*y}$.
\end{solution}

Particularly relevant for us will be two classes of linear operators which are defined next.

\begin{definition}
Let $X,Y$ be Hilbert spaces.
\begin{enumerate}
\item Let $T:X\to X$ be a bounded linear operator. Then $T$ is called \emph{self-adjoint} if $T=T^*$.
\item Let $U:X\to Y$ be a bounded linear operator. Then $U$ is called \emph{unitary} if $U$
is bijective and $U^{-1}=U^*$.
\end{enumerate}
\end{definition}

\begin{proposition}\label{prop:unitary}
A bounded linear operator $U:X\to Y$ between Hilbert spaces is unitary if and only if it is
bijective and $\|U(x)\|_Y=\|x\|_X$ for all $x\in X$. In this case, one also has
$\ip{U(x)}{U(y)}=\ip{x}{y}$ for all $x,y\in X$.
\end{proposition}
\begin{proof}
\emph{Polarization identity.} For any $x,y$ in a complex Hilbert space,
\[
\ip{x}{y}=\frac14\sum_{k=0}^3i^{-k}\,\|x+i^ky\|^2 .
\]
Indeed $\|x+i^ky\|^2=\|x\|^2+\|y\|^2+i^k\ip{x}{y}+i^{-k}\ip{y}{x}$; multiplying by $i^{-k}$
and summing over $k$, the terms with $\sum_ki^{-k}=0$ and $\sum_ki^{-2k}=\sum_k(-1)^k=0$
vanish and $4\ip{x}{y}$ remains. (In the real case:
$\ip{x}{y}=\frac14(\|x+y\|^2-\|x-y\|^2)$.)

$(\Rightarrow)$ If $U^{-1}=U^*$, then $\ip{Ux}{Uy}=\ip{x}{U^*Uy}=\ip{x}{y}$; with $y=x$
we get $\|Ux\|=\|x\|$. $(\Leftarrow)$ If $U$ preserves norms, then by polarization (and
linearity of $U$) $\ip{Ux}{Uy}=\ip{x}{y}$, i.e.\ $\ip{x}{U^*Uy}=\ip{x}{y}$ for all $x$, so
$U^*U=\mathrm{Id}_X$. Since $U$ is bijective, multiplying from the right by $U^{-1}$ gives
$U^*=U^{-1}$.
\end{proof}

\begin{example}
In quantum mechanics, the configuration (state) of a quantum system is usually described by
means of a Hilbert space $X$: each vector $x\in X$ such that $\|x\|_X=1$ corresponds to a
configuration. Then an action on the quantum configuration, e.g.\ a calculation of a
quantum computer, corresponds to a bounded linear operator $T:X\to X$. In order that this
operator be well-defined resp.\ acting in a reversible way, it needs to send a configuration
to a configuration, that is $\|T(x)\|=1$ when $\|x\|=1$, resp.\ $T$ needs to be bijective.
By Proposition~\ref{prop:unitary}, that means exactly that $T$ is unitary.
\end{example}

\begin{addition}[(Example~4.7)]
The step ``$\|Tx\|=1$ whenever $\|x\|=1$'' $\Rightarrow$ ``$\|Tx\|=\|x\|$ for all $x$''
uses linearity: for $x\ne0$, $\|Tx\|=\|x\|\,\|T(x/\|x\|)\|=\|x\|$. Moreover, physically two
unit vectors $x$ and $e^{i\theta}x$ describe the \emph{same} state (a state is a ``ray''),
which is why density matrices $\ketbra{x}{x}$ (Section~6), which do not see the phase,
are used. (``Quantum calculus'' in the original should read ``quantum mechanics''.)
\end{addition}

\begin{exercise}
Let $T:\C^n\to\C^n$ be a linear operator given by a matrix $A=(a_{ij})$. Determine a
criterion on $A$ that characterizes that $T$ is self-adjoint.
\end{exercise}
\begin{solution}
By Exercise~\ref{exo:adjmatrix} and Proposition~\ref{prop:matrix}(3), $T=T^*$ iff $A=A^*$,
i.e.\ $a_{ij}=\conj{a_{ji}}$ for all $i,j$ ($A$ is \emph{Hermitian}). In particular the
diagonal entries are real. For real matrices: $A$ symmetric.
\end{solution}

\begin{exercise}\label{exo:unitcols}
Let $T:\C^n\to\C^n$ be unitary given by a matrix $A=(a_{ij})$. For $j=1,\dots,n$ let
$b_j=(a_{ij})_{i=1}^n$ be the column vectors of $A$, and for $i=1,\dots,n$ let
$c_i=(a_{ij})_{j=1}^n$ be the row vectors of $A$. Show that the $b_j$ are mutually
orthogonal and normalized, that is, $\ip{b_j}{b_k}=0$ if $j\ne k$ and $\|b_j\|=1$ for all
$j=1,\dots,n$. Show also that the $c_i$ are mutually orthogonal and normalized, too.

In fact, these properties characterize unitaries.
\end{exercise}
\begin{solution}
Unitarity means $A^*A=I=AA^*$. Now
$(A^*A)_{jk}=\sum_i\conj{a_{ij}}a_{ik}=\ip{b_j}{b_k}$, so $A^*A=I$ says exactly
$\ip{b_j}{b_k}=\delta_{jk}$. Similarly $(AA^*)_{ik}=\sum_ja_{ij}\conj{a_{kj}}=\ip{c_k}{c_i}$,
so $AA^*=I$ says $\ip{c_k}{c_i}=\delta_{ik}$.

\emph{Converse (``characterize''):} if the columns are orthonormal, then $A^*A=I$; hence
$A$ is injective ($Ax=0\Rightarrow x=A^*Ax=0$), thus bijective as a square matrix, and
$A^{-1}=A^*$ (multiply $A^*A=I$ by $A^{-1}$ from the right). If the rows are orthonormal,
$AA^*=I$ and the same argument applies to $A^*$. So in finite dimension, orthonormal
columns $\iff$ orthonormal rows $\iff$ unitary.
\end{solution}

From now on, we will say that a matrix is self-adjoint (resp.\ unitary) whenever the
associated operator is self-adjoint (resp.\ unitary).

\begin{definition}
Let $n\in\N$. Let $A=(a_{ij})$ be an $n\times n$ matrix. It is called \emph{diagonal} if
$a_{ij}=0$ for all indices $i\ne j$. That is, $A$ has the shape
\[
A=\begin{pmatrix}a_{11}&0&0&\dots&0\\0&a_{22}&0&\dots&0\\\vdots&\vdots&\vdots&&\vdots\\
0&0&\dots&a_{n-1,n-1}&0\\0&0&\dots&0&a_{nn}\end{pmatrix}.
\]
\end{definition}

\begin{exercise}
Determine when exactly a diagonal matrix gives rise to a self-adjoint resp.\ unitary operator.
\end{exercise}
\begin{solution}
$A^*=\mathrm{diag}(\conj{a_{11}},\dots,\conj{a_{nn}})$ and
$A^*A=AA^*=\mathrm{diag}(|a_{11}|^2,\dots,|a_{nn}|^2)$. Hence a diagonal matrix is
self-adjoint iff all $a_{ii}\in\R$, and unitary iff $|a_{ii}|=1$ for all $i$
(i.e.\ $a_{ii}=e^{i\theta_i}$).
\end{solution}

\begin{theorem}[Spectral theorem]\label{thm:spectral}
Let $n\in\N$. Let $A$ be a self-adjoint or unitary $n\times n$ matrix. Then there exists a
diagonal matrix $D$ and a unitary matrix $U$, both of size $n\times n$ again, such that
\[
A=UDU^{-1}\ (=UDU^*).
\]
The entries $\lambda$ on the diagonal of $D$ are exactly the eigenvalues of $A$: there
exists $x\in\C^n$, $x\ne0$ such that $Ax=\lambda x$. The columns of $U$ are the
corresponding eigenvectors, normalized so that $\|x\|=1$. In the $i$-th column of $U$ it is
the eigenvector of the eigenvalue on the $i$-th diagonal entry of $D$. If $A$ has one or
several eigenspaces of dimension $\geqslant2$, then one takes several orthogonal normalized
vectors in $U$. If $A$ is self-adjoint, its eigenvalues are real; if $A$ is unitary, they
have modulus $1$.
\end{theorem}

\begin{addition}[(proof of the spectral theorem)]
Both self-adjoint and unitary matrices are \emph{normal}: $AA^*=A^*A$. We prove the
theorem for normal $A$, by induction on $n$; $n=1$ is trivial.

\emph{Step 1: for normal $A$, $Ax=\lambda x\iff A^*x=\conj\lambda x$.} $B=A-\lambda I$ is
normal too, and $\|Bx\|^2=\ip{x}{B^*Bx}=\ip{x}{BB^*x}=\|B^*x\|^2$. As $B^*=A^*-\conj\lambda I$,
the claim follows.

\emph{Step 2.} By the fundamental theorem of algebra, the characteristic polynomial
$\det(A-\lambda I)$ has a root $\lambda_1$, i.e.\ $A$ has an eigenvector $u_1$ with
$\|u_1\|=1$. Let $V=u_1^\perp=\{v:\ip{u_1}{v}=0\}$, of dimension $n-1$. For $v\in V$:
$\ip{u_1}{Av}=\ip{A^*u_1}{v}=\ip{\conj{\lambda_1}u_1}{v}=\lambda_1\ip{u_1}{v}=0$ and
$\ip{u_1}{A^*v}=\ip{Au_1}{v}=\conj{\lambda_1}\ip{u_1}{v}=0$. So $A$ and $A^*$ map $V$ into
$V$; the restriction $A|_V$ is normal on the Hilbert space $V$ (with adjoint $A^*|_V$). By
induction $V$ has an orthonormal basis $u_2,\dots,u_n$ of eigenvectors of $A$.

\emph{Step 3.} Let $U$ be the matrix with columns $u_1,\dots,u_n$ and
$D=\mathrm{diag}(\lambda_1,\dots,\lambda_n)$. $U$ is unitary by
Exercise~\ref{exo:unitcols}, and $AU=UD$ column by column, hence $A=UDU^{-1}$.
Conversely, if $A=UDU^{-1}$ then $Au_i=d_{ii}u_i$ and
$\det(A-\lambda I)=\det(D-\lambda I)=\prod_i(d_{ii}-\lambda)$, so the $d_{ii}$ are all the
eigenvalues, with multiplicity.

\emph{Eigenvalues.} If $A=A^*$ and $Ax=\lambda x$, $\|x\|=1$:
$\lambda=\ip{x}{Ax}=\ip{Ax}{x}=\conj\lambda$. If $A$ is unitary:
$|\lambda|=\|Ax\|=\|x\|=1$. (If $A$ is real symmetric, all of this can be done in $\R^n$,
and $U$ can be chosen real orthogonal.)
\end{addition}

\begin{exercise}
Consider the matrix
\[
A=\begin{pmatrix}1&-\frac23&-\frac23\\-\frac23&\frac23&0\\-\frac23&0&\frac43\end{pmatrix}.
\]
\begin{enumerate}
\item Is this matrix self-adjoint? Is it unitary?
\item The following are eigenvectors of $A$: $x=(1,-2,2)$, $y=(2,-1,-2)$, $z=(2,2,1)$.
With this information at hand, try to find matrices $D$ and $U$ as in
Theorem~\ref{thm:spectral} such that $A=UDU^{-1}$.
\end{enumerate}
\end{exercise}
\begin{solution}
1. $A$ is real and symmetric, hence $A=A^*$: \emph{self-adjoint}. It is \emph{not unitary}:
its first column has norm$^2$ $1+\frac49+\frac49=\frac{17}9\ne1$ (Exercise~\ref{exo:unitcols});
equivalently $A^*A=A^2\ne I$. (It is not even invertible, see below.)

2. Compute $Ax=(1+\frac43-\frac43,\,-\frac23-\frac43,\,-\frac23+\frac83)=(1,-2,2)=1\cdot x$,
$Ay=(2+\frac23+\frac43,\,-\frac43-\frac23,\,-\frac43-\frac83)=(4,-2,-4)=2\cdot y$,
$Az=(2-\frac43-\frac23,\,-\frac43+\frac43,\,-\frac43+\frac43)=(0,0,0)=0\cdot z$.
The eigenvalues are $1,2,0$ (check: $\tr A=1+\frac23+\frac43=3=1+2+0$). The eigenvectors
are pairwise orthogonal ($\ip xy=2+2-4=0$, $\ip xz=2-4+2=0$, $\ip yz=4-2-2=0$, as it must
be for distinct eigenvalues of a self-adjoint matrix) and each has norm $3$. Hence
\[
U=\frac13\begin{pmatrix}1&2&2\\-2&-1&2\\2&-2&1\end{pmatrix},\qquad
D=\begin{pmatrix}1&0&0\\0&2&0\\0&0&0\end{pmatrix},\qquad U^{-1}=U^*=U^T,
\]
and $A=UDU^{-1}$. Since $0$ is an eigenvalue, $A$ is not invertible.
\end{solution}

\begin{definition}\label{def:HS}
Let $n\in\N$ and denote $M_n(\C)$ the space of $n\times n$ matrices (with complex entries).
The \emph{trace} $\tr:M_n(\C)\to\C$ is the linear functional $\tr((a_{ij})_{ij})=\sum_{i=1}^na_{ii}$.
It defines a scalar product on $M_n(\C)$ by
\[
\ip{A}{B}=\tr(A^*B).
\]
Here, $A^*$ is the adjoint matrix of $A$. In this way, $M_n(\C)$ becomes a Hilbert space;
the above scalar product and its norm are called \emph{Hilbert--Schmidt scalar product} and
\emph{Hilbert--Schmidt norm}.
\end{definition}

\begin{addition}[(why this is a scalar product)]
$\tr(A^*B)=\sum_i\sum_j(A^*)_{ij}b_{ji}=\sum_{i,j}\conj{a_{ji}}b_{ji}$, i.e.\ the
Hilbert--Schmidt product is the standard scalar product of $\C^{n^2}$ applied to the
entries. Hence it is a scalar product, $M_n(\C)$ is a Hilbert space (Example~\ref{ex:Cn}),
and $\|A\|_{HS}^2=\sum_{i,j}|a_{ij}|^2$. Useful facts: $\tr(AB)=\tr(BA)$ (both equal
$\sum_{i,j}a_{ij}b_{ji}$), and, by Cauchy--Schwarz, $\|AB\|_{HS}\le\|A\|_{HS}\|B\|_{HS}$ and
$\|Ax\|\le\|A\|_{HS}\|x\|$.
\end{addition}

\begin{definition}
Let $n\in\N$. Let $A$ be an $n\times n$ matrix. For any $k\geqslant1$, the $k$-th power of
$A$ is defined as $A^k=A\cdot A\cdots A$ with $k$ factors. For $k=0$, $A^0=1_n$ where $1_n$
is the diagonal matrix with entries $1$ everywhere on the diagonal.

We define the \emph{matrix exponential}
\[
e^A=\sum_{k=0}^\infty\frac1{k!}A^k .
\]
Similarly for $z\in\C$, we define the matrix exponential
\[
e^{zA}=\sum_{k=0}^\infty\frac{z^k}{k!}A^k .
\]
The convergence of the series is e.g.\ with respect to the Hilbert--Schmidt norm.
\end{definition}

\begin{addition}[(convergence)]
By submultiplicativity, $\|A^k\|_{HS}\le\|A\|_{HS}^k$ for $k\ge1$, so
$\sum_k\frac{|z|^k}{k!}\|A^k\|_{HS}\le\sqrt n+e^{|z|\|A\|_{HS}}<\infty$. An absolutely
convergent series in the complete space $M_n(\C)$ converges (the partial sums form a Cauchy
sequence). Since all norms on the finite-dimensional space $M_n(\C)$ are equivalent, the
limit is the same for any norm, e.g.\ the operator norm of Section~5. Also: if
$A=UDU^{-1}$ with $D=\mathrm{diag}(\lambda_i)$, then $A^k=UD^kU^{-1}$ and
$e^{zA}=U\,\mathrm{diag}(e^{z\lambda_i})\,U^{-1}$.
\end{addition}

A particularly important example of matrices in quantum mechanics are the following
$2\times2$ matrices.

\begin{definition}
Consider the following $2\times2$ matrices:
\[
\sigma_x=\begin{pmatrix}0&1\\1&0\end{pmatrix},\quad
\sigma_y=\begin{pmatrix}0&-i\\i&0\end{pmatrix},\quad
\sigma_z=\begin{pmatrix}1&0\\0&-1\end{pmatrix},\quad
1=\begin{pmatrix}1&0\\0&1\end{pmatrix}.
\]
The matrices $\sigma_x,\sigma_y,\sigma_z$ are called \emph{Pauli matrices}.
\end{definition}

\begin{exercise}
\begin{enumerate}
\item Convince yourself that $\sigma_x,\sigma_y,\sigma_z$ are all both self-adjoint and unitary.
\item Show that $\sigma_x$ and $\sigma_y$ are orthogonal in the Hilbert--Schmidt scalar product.
\end{enumerate}
In fact, the matrices $1,\sigma_x,\sigma_y,\sigma_z$ are all mutually orthogonal and form
a basis of the $4$-dimensional $\C$-vector space $M_2(\C)$.
\end{exercise}
\begin{solution}
1. Each matrix equals its conjugate transpose (for $\sigma_y$:
$\conj{\begin{psmallmatrix}0&-i\\i&0\end{psmallmatrix}}^{T}=\begin{psmallmatrix}0&-i\\i&0\end{psmallmatrix}$).
Moreover $\sigma_x^2=\sigma_y^2=\sigma_z^2=1$ (direct computation), so
$\sigma^{-1}=\sigma=\sigma^*$: unitary.

2. $\sigma_x\sigma_y=\begin{psmallmatrix}i&0\\0&-i\end{psmallmatrix}=i\sigma_z$, so
$\ip{\sigma_x}{\sigma_y}=\tr(\sigma_x^*\sigma_y)=\tr(i\sigma_z)=0$.

\emph{The remark.} Similarly $\sigma_y\sigma_z=i\sigma_x$, $\sigma_z\sigma_x=i\sigma_y$,
and all Pauli matrices are traceless, so $\ip{1}{\sigma_j}=\tr\sigma_j=0$ and
$\ip{\sigma_j}{\sigma_k}=\tr(\sigma_j\sigma_k)=0$ for $j\ne k$; also
$\ip{\sigma_j}{\sigma_j}=\tr1=2=\ip11$. Four mutually orthogonal non-zero vectors are
linearly independent, hence a basis of the 4-dimensional space $M_2(\C)$. Explicitly,
\[
\begin{pmatrix}a&b\\c&d\end{pmatrix}=\frac{a+d}2\,1+\frac{b+c}2\,\sigma_x+\frac{i(b-c)}2\,\sigma_y+\frac{a-d}2\,\sigma_z .\qedhere
\]
\end{solution}

\begin{exercise}\label{exo:expsx}
\begin{enumerate}
\item Calculate $\sigma_x^2$.
\item Deduce $\sigma_x^k$ for all $k\geqslant0$.
\item Deduce the matrix exponential $\exp(z\sigma_x)$ for $z\in\C$.
\end{enumerate}
\end{exercise}
\begin{solution}
1. $\sigma_x^2=\begin{psmallmatrix}0&1\\1&0\end{psmallmatrix}\begin{psmallmatrix}0&1\\1&0\end{psmallmatrix}=1$.

2. $\sigma_x^{2m}=(\sigma_x^2)^m=1$ and $\sigma_x^{2m+1}=\sigma_x$.

3. Splitting the series into even and odd terms,
\[
e^{z\sigma_x}=\sum_{m\ge0}\frac{z^{2m}}{(2m)!}1+\sum_{m\ge0}\frac{z^{2m+1}}{(2m+1)!}\sigma_x
=\cosh(z)\,1+\sinh(z)\,\sigma_x=\begin{pmatrix}\cosh z&\sinh z\\\sinh z&\cosh z\end{pmatrix}.
\]
In particular, for $z=-i\theta$: $e^{-i\theta\sigma_x}=\cos\theta\,1-i\sin\theta\,\sigma_x$
(used in Example~8.4). The same holds for $\sigma_y,\sigma_z$ and in fact for
$\vec n\cdot\vec\sigma$ with $\vec n\in\R^3$ a unit vector, since $(\vec n\cdot\vec\sigma)^2=1$.
\end{solution}
%======================================================================
\section{Operator exponential and Trotter--Lie product formula}

We recall that in Definition~\ref{def:bdd}, we have defined the notion of a bounded linear
operator between Hilbert spaces. That is, if $X$ is a Hilbert space (more generally a
complete normed space), then a bounded linear mapping $A:X\to X$ satisfies by definition
$A(x+\lambda y)=A(x)+\lambda A(y)$ for $x,y\in X$ and $\lambda\in\C$ and for some constant
$C>0$, we have $\|A(x)\|_X\le C\|x\|_X$.

\begin{definition}
If $A:X\to X$ is a bounded linear operator as above, then the least constant $C\geqslant0$
such that $\|A(x)\|\le C\|x\|$ holds for all $x\in X$ is denoted by $\|A\|$. Equivalently,
$\|A\|=\sup_{\|x\|\le1}\|A(x)\|$ (the \emph{operator norm}).
\end{definition}

\begin{remark}
The space $B(X)$ of bounded linear operators from $X$ to $X$ is a normed vector space with
norm $\|A\|$ as in Definition~5.1 above.
\end{remark}

\begin{addition}[(proof of Remark 5.2, and completeness of $B(X)$)]
\emph{Least constant exists:} with $s=\sup_{\|x\|\le1}\|Ax\|$, for $x\ne0$ we have
$\|Ax\|=\|x\|\,\|A(x/\|x\|)\|\le s\|x\|$, and no $C<s$ works. \emph{Norm axioms:}
$\|A\|=0\Rightarrow Ax=0$ for all $x$; $\|\lambda A\|=|\lambda|\|A\|$ is clear;
$\|(A+B)x\|\le\|Ax\|+\|Bx\|\le(\|A\|+\|B\|)\|x\|$.

\emph{$B(X)$ is complete} (a Banach space): let $(A_n)$ be Cauchy in $B(X)$. For each $x$,
$\|A_nx-A_mx\|\le\|A_n-A_m\|\|x\|$, so $(A_nx)$ is Cauchy in $X$ and converges to some
$Ax$. $A$ is linear (limits of linear expressions), and given $\epsilon>0$ and $n_0$ with
$\|A_n-A_m\|\le\epsilon$ for $n,m\ge n_0$, letting $m\to\infty$ in
$\|A_nx-A_mx\|\le\epsilon\|x\|$ gives $\|(A_n-A)x\|\le\epsilon\|x\|$. Hence $A\in B(X)$ and
$\|A_n-A\|\le\epsilon$ for $n\ge n_0$.

\emph{The ``result from Banach space theory'' used in Lemma~5.5:} in a complete normed space,
every absolutely convergent series $\sum_n a_n$ ($\sum_n\|a_n\|<\infty$) converges, because
$\|\sum_{n=N}^Ma_n\|\le\sum_{n=N}^M\|a_n\|\to0$, so the partial sums are Cauchy.
\end{addition}

\begin{lemma}\label{lem:comp}
Let $A,B:X\to X$ be two bounded linear operators on a Hilbert space $X$. Then
$A\circ B:X\to X$ is again a bounded linear operator and $\|A\circ B\|\le\|A\|\,\|B\|$.
\end{lemma}
\begin{proof}
We have
$A\circ B(x+\lambda y)=A(B(x+\lambda y))=A(B(x)+\lambda B(y))=A(B(x))+\lambda A(B(y))
=A\circ B(x)+\lambda A\circ B(y)$.
Thus, $A\circ B$ is again linear. Moreover,
\[
\|A\circ B(x)\|=\|A(B(x))\|\le\|A\|\,\|B(x)\|\le\|A\|\,\|B\|\,\|x\|.
\]
Thus $A\circ B$ is bounded and $\|A\circ B\|\le\|A\|\,\|B\|$.
\end{proof}

\begin{definition}\label{def:opexp}
Let $X$ and $A$ be as above. Let $t\in\R$ (or even $t\in\C$). Then we define the
\emph{operator exponential} $e^{tA}$ by
\[
e^{tA}=\sum_{n=0}^\infty\frac{(tA)^n}{n!}.
\]
\end{definition}

\begin{correction}[(Def.~5.4)]
The original defines $e^{tA}$ only for $t\ge0$. However, the proof of Proposition~5.7 uses
$e^{hA}$ for $h\ne0$ of either sign (to compute a two-sided derivative), and the inverse of
$e^{tA}$ is $e^{-tA}$. Nothing in the definition or in Lemmas 5.5--5.6 requires $t\ge0$, so
we define $e^{tA}$ for all real (or complex) $t$.
\end{correction}

Note that by Lemma~\ref{lem:comp}, we have indeed $(tA)^n\in B(X)$.

\begin{lemma}\label{lem:abs}
The series defining $e^{tA}$ above converges absolutely, and by a result from Banach space
theory (see the Complement after Remark 5.2) therefore also simply.
Moreover $\|e^{tA}\|\le e^{|t|\,\|A\|}$.
\end{lemma}
\begin{proof}
Indeed, by Lemma~\ref{lem:comp} applied successively, we have
\[
\|A^n\|=\|A^{n-1}\circ A\|\le\|A^{n-1}\|\,\|A\|\le\dots\le\|A\|\,\|A\|\cdots\|A\|=\|A\|^n.
\]
Thus, $\|(tA)^n\|=\|t^nA^n\|=|t|^n\|A^n\|\le|t|^n\|A\|^n$ and
\[
\sum_{n=0}^\infty\Big\|\frac{(tA)^n}{n!}\Big\|\le\sum_{n=0}^\infty\frac{|t|^n\|A\|^n}{n!}=e^{|t|\,\|A\|}<\infty .\qedhere
\]
\end{proof}

As for the scalar valued exponential, using Cauchy summation of double sums, one can prove
the following lemma.

\begin{lemma}\label{lem:explaw}
The operator exponential satisfies the exponential law
\[
e^{(t+s)A}=e^{tA}e^{sA}\qquad(t,s\in\R).
\]
In particular $e^{tA}$ is invertible with $(e^{tA})^{-1}=e^{-tA}$, and $(e^{A/k})^k=e^A$ for $k\in\N$.
\end{lemma}

\begin{addition}[(proof of Lemma~5.6)]
\emph{Cauchy product in $B(X)$.} Let $a_n,b_n\in B(X)$ with $\sum\|a_n\|<\infty$,
$\sum\|b_n\|<\infty$ and put $c_n=\sum_{k=0}^na_kb_{n-k}$. Then
$\sum_nc_n=\big(\sum_ka_k\big)\big(\sum_lb_l\big)$. Indeed,
\begin{align*}
\Big\|\sum_{n\le N}c_n-\sum_{k\le N}a_k\sum_{l\le N}b_l\Big\|
&=\Big\|\sum_{\substack{k,l\le N\\k+l>N}}a_kb_l\Big\|\\
&\le\Big(\sum_k\|a_k\|\Big)\sum_{l>N/2}\|b_l\|+\Big(\sum_{k>N/2}\|a_k\|\Big)\sum_l\|b_l\|\xrightarrow[N\to\infty]{}0,
\end{align*}
because $k+l>N$ forces $k>N/2$ or $l>N/2$, and tails of convergent series tend to $0$.

\emph{Exponential law.} Apply this with $a_k=\frac{(tA)^k}{k!}$, $b_l=\frac{(sA)^l}{l!}$
(absolutely convergent by Lemma~\ref{lem:abs}). Since powers of the \emph{same} operator
$A$ commute, $a_kb_{n-k}=\frac{t^ks^{n-k}}{k!(n-k)!}A^n$, and the binomial theorem gives
\[
c_n=\frac{A^n}{n!}\sum_{k=0}^n\binom nkt^ks^{n-k}=\frac{(t+s)^nA^n}{n!},\qquad
e^{tA}e^{sA}=\sum_nc_n=e^{(t+s)A}.
\]
With $s=-t$: $e^{tA}e^{-tA}=e^{0A}=I=e^{-tA}e^{tA}$. Induction gives $(e^{A/k})^k=e^{A}$.

\emph{Remark.} The same proof shows $e^{A+B}=e^Ae^B$ whenever $AB=BA$ (the binomial
formula $(A+B)^n=\sum\binom nkA^kB^{n-k}$ requires commutation). Example~5.8 shows that this
fails in general.
\end{addition}

The operator exponential can be used to solve abstractly an initial value differential
equation. Namely, we let $X$ be a Hilbert space and $A:X\to X$ a bounded linear operator.
Let $u_0\in X$ be an initial value. Consider the differential equation
\begin{equation}\label{eq:ivp}
\begin{cases}\partial_tu(t)=Au(t)&(t\ge0)\\u(0)=u_0\end{cases},
\end{equation}
where $u:[0,\infty)\to X$ is the unknown function to find.

\begin{proposition}\label{prop:ivp}
The solution of \eqref{eq:ivp} is given by $u(t)=e^{tA}u_0$ for $t\geqslant0$, and it is the
only solution.
\end{proposition}

We shall prove that the above $u(t)$ solves indeed the differential equation, and then
that it is the only solution.

\begin{proof}[Proof (existence)]
Let us show that $u$ is differentiable and calculate at the same time its derivative. We
have for $h\ne0$ and $t\ge0$ (with $t+h\ge0$), according to Lemma~\ref{lem:explaw},
\[
\frac1h\big(u(t+h)-u(t)\big)=\frac1h\big(e^{(t+h)A}u_0-e^{tA}u_0\big)=\frac1h\big(e^{hA}e^{tA}u_0-e^{tA}u_0\big).
\]
Write $u_1=e^{tA}u_0$. Then the above calculation continues
\[
=\frac1h\big(e^{hA}u_1-u_1\big)=\frac1h\Big(\sum_{n=0}^\infty\frac{(hA)^n}{n!}u_1-u_1\Big)
=\frac1h\sum_{n=1}^\infty\frac{(hA)^n}{n!}u_1=\sum_{n=1}^\infty A\frac{(hA)^{n-1}}{n!}u_1 .
\]
We estimate the norm of the difference:
\begin{align*}
\Big\|\frac1h\big(u(t+h)-u(t)\big)-Ae^{tA}u_0\Big\|
&=\Big\|\sum_{n=1}^\infty A\frac{(hA)^{n-1}}{n!}u_1-Au_1\Big\|
=\Big\|\sum_{n=2}^\infty A\frac{(hA)^{n-1}}{n!}u_1\Big\|\\
&\le\sum_{n=2}^\infty\frac{|h|^{n-1}\|A\|^{n}}{n!}\|u_1\|
\le|h|\,\|A\|^2\sum_{n=2}^\infty\frac{|h|^{n-2}\|A\|^{n-2}}{(n-2)!}\|u_1\|\\
&=|h|\,\|A\|^2e^{|h|\,\|A\|}\|u_1\|\to0\qquad(h\to0).
\end{align*}
Thus, $u(t)$ is differentiable in $t$ and $\partial_tu(t)=Ae^{tA}u_0=Au(t)$. Moreover,
$u(0)=e^{0A}u_0=\big(I+\sum_{n\ge1}\frac{0^n}{n!}A^n\big)u_0=u_0$.
\end{proof}

\begin{correction}[(proof of Prop.~5.7)]
(a) The reference ``Lemma ??'' is broken; it is Lemma~5.6 (exponential law).
(b) In the second line of the estimate, the original writes
$\big\|\sum_{n=2}^\infty\frac{(hA)^{n-1}}{n!}u_1\big\|$: the factor $A$ in front was lost.
Consequently the original final bound $|h|\,\|A\|\,e^{|h|\|A\|}\|u_1\|$ should be
$|h|\,\|A\|^2e^{|h|\|A\|}\|u_1\|$, as above. Since it still tends to $0$, the conclusion
is unaffected. (c) For $h<0$ one needs $e^{hA}$ with negative exponent, see the correction
to Definition~5.4; at $t=0$ only $h>0$ is allowed (right derivative).
\end{correction}

\begin{proof}[Proof (uniqueness; ``with a little more work'')]
The estimate above, with $u_1$ replaced by an arbitrary vector of norm $\le1$, shows
$\big\|\frac1h(e^{(t+h)A}-e^{tA})-Ae^{tA}\big\|\le|h|\|A\|^2e^{|h|\|A\|}\|e^{tA}\|\to0$,
i.e.\ $t\mapsto e^{tA}$ is differentiable in operator norm with derivative $Ae^{tA}=e^{tA}A$;
in particular it is continuous. Let now $v$ be any solution of \eqref{eq:ivp} and
$w(t)=e^{-tA}v(t)$. Then
\begin{align*}
\frac{w(t+h)-w(t)}h&=e^{-(t+h)A}\frac{v(t+h)-v(t)}h+\frac{e^{-(t+h)A}-e^{-tA}}h\,v(t)\\
&\xrightarrow[h\to0]{}e^{-tA}Av(t)-Ae^{-tA}v(t)=0,
\end{align*}
using $\|e^{-(t+h)A}\|\le e^{(|t|+1)\|A\|}$ for $|h|\le1$ and the continuity of
$s\mapsto e^{-sA}$. So $w'\equiv0$ on $[0,\infty)$. For every $y\in X$, the scalar
function $t\mapsto\ip{y}{w(t)}$ has zero derivative (real and imaginary parts separately),
hence is constant by the mean value theorem: $\ip{y}{w(t)-w(0)}=0$ for all $y$, so
$w(t)=w(0)=v(0)=u_0$. Thus $v(t)=e^{tA}w(t)=e^{tA}u_0=u(t)$.
\end{proof}

In practice, the operator exponential is often of the form $e^{t(A+B)}$ where $A$ and $B$ are
two operators of different nature. One can often calculate $e^{tA}$ and $e^{tB}$
separately. The following example shows that this does not automatically give the full
operator exponential $e^{t(A+B)}$.

\begin{example}\label{ex:noncomm}
Let $X=\C^2$ and
\[
A=\begin{pmatrix}0&1\\0&0\end{pmatrix},\qquad B=\begin{pmatrix}0&0\\1&0\end{pmatrix}.
\]
We can calculate powers of $A$:
$A^0=\begin{psmallmatrix}1&0\\0&1\end{psmallmatrix}=:I$, $A^1=A$ and
$A^2=\begin{psmallmatrix}0&1\\0&0\end{psmallmatrix}\begin{psmallmatrix}0&1\\0&0\end{psmallmatrix}=0$.
Thus for any $n\ge2$, $A^n=0$. Therefore the operator exponential is
\[
e^{tA}=I+tA=\begin{pmatrix}1&t\\0&1\end{pmatrix}.
\]
In a similar manner, $B^0=I$, $B^1=B$, $B^n=0$ for $n\ge2$, and therefore
$e^{tB}=I+tB=\begin{psmallmatrix}1&0\\t&1\end{psmallmatrix}$. The product of the two
operator exponentials becomes
\[
e^{tA}e^{tB}=\begin{pmatrix}1&t\\0&1\end{pmatrix}\begin{pmatrix}1&0\\t&1\end{pmatrix}
=\begin{pmatrix}1+t^2&t\\t&1\end{pmatrix}.
\]
On the other hand, $A+B=\begin{psmallmatrix}0&1\\1&0\end{psmallmatrix}\,(=\sigma_x)$ and
$(A+B)^2=I$. Therefore $(A+B)^n=A+B$ for $n$ odd and $(A+B)^n=I$ for $n$ even. We deduce
the operator exponential of $A+B$ (compare Exercise~\ref{exo:expsx}):
\[
e^{t(A+B)}=\sum_{n=0}^\infty\frac{t^n}{n!}(A+B)^n=\sum_{n=0}^\infty\frac{t^{2n}}{(2n)!}I
+\sum_{n=0}^\infty\frac{t^{2n+1}}{(2n+1)!}(A+B)=\begin{pmatrix}\cosh t&\sinh t\\\sinh t&\cosh t\end{pmatrix}.
\]
We see in this example that in general, $e^{tA}e^{tB}\ne e^{t(A+B)}$.
\end{example}

\begin{addition}[(why: non-commutativity)]
Here $AB=\begin{psmallmatrix}1&0\\0&0\end{psmallmatrix}\ne\begin{psmallmatrix}0&0\\0&1\end{psmallmatrix}=BA$,
with commutator $[A,B]=AB-BA=\sigma_z$. Expanding to second order,
$e^{tA}e^{tB}-e^{t(A+B)}=\begin{psmallmatrix}t^2/2&0\\0&-t^2/2\end{psmallmatrix}+O(t^3)
=\frac{t^2}2[A,B]+O(t^3)$, in accordance with the Baker--Campbell--Hausdorff formula
$e^{tA}e^{tB}=e^{t(A+B)+\frac{t^2}2[A,B]+O(t^3)}$.
\end{addition}

In order to obtain a correct dependency of $e^{t(A+B)}$ in terms of $e^{tA}$ and $e^{tB}$,
there is the following theorem.

\begin{theorem}[Trotter--Lie product formula]\label{thm:trotter}
Let $X$ be a Hilbert space and $A,B$ be bounded linear operators over $X$. Then for any
$t\geqslant0$, we have the Trotter--Lie formula
\[
e^{t(A+B)}=\lim_{n\to\infty}\big(e^{\frac tnA}e^{\frac tnB}\big)^n
\]
(limit in operator norm). More precisely,
$\big\|e^{t(A+B)}-(e^{\frac tnA}e^{\frac tnB})^n\big\|\le\frac2n\,e^{2t(\|A\|+\|B\|)}$.
\end{theorem}
\begin{proof}
Upon renaming $tA$ into $A$ and $tB$ into $B$, we can assume that $t=1$. We write $k$
instead of $n$. Define $C=e^{(A+B)/k}$ and $D=e^{A/k}e^{B/k}$. By Lemma~\ref{lem:explaw},
$C^k=e^{A+B}$, and $D^k$ is the Trotter product, so we must show $\|C^k-D^k\|\to0$.
We have the following estimates according to the proof of Lemma~\ref{lem:abs} and
Lemma~\ref{lem:comp}:
\[
\|C\|,\ \|D\|\le\exp\Big[\frac{\|A\|+\|B\|}k\Big]=\exp(\|A\|+\|B\|)^{1/k}.
\]
Applying the Cauchy product formula (Complement after Lemma~5.6) to $D$ yields
\[
D=e^{A/k}e^{B/k}=\sum_{i=0}^\infty\frac{(A/k)^i}{i!}\cdot\sum_{j=0}^\infty\frac{(B/k)^j}{j!}
=\sum_{m=0}^\infty k^{-m}\sum_{i=0}^m\frac{A^i}{i!}\cdot\frac{B^{m-i}}{(m-i)!}.
\]
The terms with $m=0$ ($I$) and $m=1$ ($\frac{A+B}k$) coincide in $C$ and $D$. This permits
us to bound the norm of $C-D$ by
\begin{align*}
\|C-D\|&=\Big\|\sum_{i=0}^\infty\frac{([A+B]/k)^i}{i!}-\sum_{m=0}^\infty k^{-m}\sum_{i=0}^m\frac{A^i}{i!}\cdot\frac{B^{m-i}}{(m-i)!}\Big\|\\
&=\Big\|\sum_{i=2}^\infty k^{-i}\frac{(A+B)^i}{i!}-\sum_{m=2}^\infty k^{-m}\sum_{i=0}^m\frac{A^i}{i!}\cdot\frac{B^{m-i}}{(m-i)!}\Big\|\\
&\le\frac1{k^2}\Big[\exp(\|A\|+\|B\|)+\sum_{m=2}^\infty\sum_{i=0}^m\frac{\|A\|^i}{i!}\cdot\frac{\|B\|^{m-i}}{(m-i)!}\Big]\\
&=\frac1{k^2}\Big[\exp(\|A\|+\|B\|)+\sum_{m=2}^\infty\frac1{m!}\sum_{i=0}^m\binom mi\|A\|^i\|B\|^{m-i}\Big]\\
&=\frac1{k^2}\Big[\exp(\|A\|+\|B\|)+\sum_{m=2}^\infty\frac{(\|A\|+\|B\|)^m}{m!}\Big]
\le\frac2{k^2}\exp(\|A\|+\|B\|),
\end{align*}
where we used $k^{-i}\le k^{-2}$ for $i\ge2$ and $\|(A+B)^i\|\le(\|A\|+\|B\|)^i$.
By the telescoping identity
$C^k-D^k=\sum_{m=0}^{k-1}C^m(C-D)D^{k-m-1}$
(the sum equals $\sum_m(C^{m+1}D^{k-m-1}-C^mD^{k-m})$), we obtain
\[
\|C^k-D^k\|=\Big\|\sum_{m=0}^{k-1}C^m(C-D)D^{k-m-1}\Big\|\le\exp(\|A\|+\|B\|)\cdot k\cdot\|C-D\|,
\]
where we used that
$\|C\|^m\cdot\|D\|^{k-m-1}\le\exp(\|A\|+\|B\|)^{(k-1)/k}\le\exp(\|A\|+\|B\|)$. Together
with the estimate for $\|C-D\|$, this yields
\[
\|C^k-D^k\|\le\frac{2\exp(2\|A\|+2\|B\|)}k,
\]
which upon letting $k\to\infty$ proves the theorem.
\end{proof}

\begin{addition}[(numerical illustration and relevance)]
For Example~\ref{ex:noncomm} with $t=1$, one has
$e^{A/n}e^{B/n}=\begin{psmallmatrix}1+1/n^2&1/n\\1/n&1\end{psmallmatrix}$ and
$e^{A+B}=\begin{psmallmatrix}\cosh1&\sinh1\\\sinh1&\cosh1\end{psmallmatrix}
\approx\begin{psmallmatrix}1.543081&1.175201\\1.175201&1.543081\end{psmallmatrix}$:
\begin{center}\footnotesize\renewcommand{\arraystretch}{1.1}
\begin{tabular}{rccc}
\toprule
$n$ & $(e^{A/n}e^{B/n})^n$ & error $\|\cdot-e^{A+B}\|$ & $n\times$ error\\
\midrule
1 & $\begin{pmatrix}2&1\\1&1\end{pmatrix}$ & 0.5729 & 0.573\\
2 & $\begin{pmatrix}1.8125&1.125\\1.125&1.25\end{pmatrix}$ & 0.2975 & 0.595\\
10 & $\begin{pmatrix}1.601246&1.173094\\1.173094&1.483937\end{pmatrix}$ & 0.0592 & 0.592\\
100 & $\begin{pmatrix}1.548952&1.175180\\1.175180&1.537200\end{pmatrix}$ & 0.00588 & 0.588\\
1000 & $\begin{pmatrix}1.543668&1.175201\\1.175201&1.542493\end{pmatrix}$ & 0.000588 & 0.588\\
\bottomrule
\end{tabular}
\end{center}
The error decays exactly like $1/n$, as predicted by the theorem. In quantum computing, this
is the basis of \emph{Trotterization}: to simulate $e^{-itH}$ for a Hamiltonian
$H=H_1+H_2$ whose parts are easy to implement as gates, one applies
$(e^{-itH_1/n}e^{-itH_2/n})^n$; the theorem guarantees an error $O(t^2/n)$ (more precisely,
from the second-order term, $\approx\frac{t^2}{2n}\|[H_1,H_2]\|$).
\end{addition}
%======================================================================
\section{Density matrix}

\begin{definition}\label{def:psd}
Let $n\in\N$. A matrix $A=(a_{ij})_{ij}\in M_n(\C)$ is \emph{positive semi-definite} if
\[
\ip{x}{A(x)}=\sum_{i,j=1}^nx_j\conj{x_i}a_{ij}\geqslant0
\]
for all $x\in\C^n$. Moreover, $A$ is called \emph{positive definite} if
\[
\ip{x}{A(x)}=\sum_{i,j=1}^nx_j\conj{x_i}a_{ij}>0
\]
for all $x\in\C^n$, $x\ne0$.
\end{definition}

\begin{addedthm}{Lemma 6.1a}
Let $A\in M_n(\C)$. The following are equivalent:
(i) $A$ is positive semi-definite;
(ii) $A$ is self-adjoint and all its eigenvalues are $\ge0$;
(iii) $A=\sum_k\lambda_k\ketbra{u_k}{u_k}$ for an orthonormal basis $(u_k)$ and $\lambda_k\ge0$.
Similarly, $A$ is positive definite iff it is self-adjoint with all eigenvalues $>0$.
\end{addedthm}
\begin{proof}
(i)$\Rightarrow$(ii): $\ip{x}{Ax}\in\R$ for all $x$, so the sesquilinear form
$s(x,y)=\ip{x}{(A-A^*)y}$ satisfies $s(x,x)=\ip{x}{Ax}-\conj{\ip{x}{Ax}}=0$ for all $x$.
The polarization identity (proof of Proposition~\ref{prop:unitary}) holds for every
sesquilinear form: $s(x,y)=\frac14\sum_{k=0}^3i^{-k}s(x+i^ky,x+i^ky)=0$. Hence $A=A^*$.
If $Ax=\lambda x$, $\|x\|=1$, then $\lambda=\ip{x}{Ax}\ge0$.
(ii)$\Rightarrow$(iii) is Theorem~\ref{thm:spectral}: $A=UDU^*=\sum_k\lambda_k\ketbra{u_k}{u_k}$
with $u_k$ the columns of $U$.
(iii)$\Rightarrow$(i): $\ip{x}{Ax}=\sum_k\lambda_k|\ip{u_k}{x}|^2\ge0$.
The positive definite case is identical, using $\sum_k\lambda_k|\ip{u_k}{x}|^2>0$ for
$x\ne0$ iff all $\lambda_k>0$ (Parseval). Note that over $\R$ (i) does \emph{not} imply
symmetry, e.g.\ $\begin{psmallmatrix}0&1\\-1&0\end{psmallmatrix}$; over $\C$ it does.
\end{proof}

Let $(\Omega,\PP)$ be a probability space. We recall from Definition~\ref{def:expect} that
the expectation $\E$ has the properties that it is linear on the space of bounded functions
$\Omega\to\C$, that $\E(1)=1$ and that $\E(X)\ge0$ whenever $X\ge0$. Let now $n\in\N$. Note
that the space $M_n(\C)$ is linear, contains an element $1_n$ (that is, the diagonal matrix
with $1$ entries on the diagonal), and carries the positivity notion ``positive
semi-definite'' from Definition~\ref{def:psd}. We are thus able to define a generalized
version of expectation on the space $M_n(\C)$.

\begin{definition}
Let $\Phi:M_n(\C)\to\C$ be a function. Then $\Phi$ is called a \emph{state} if
\begin{enumerate}
\item $\Phi$ is linear,
\item $\Phi(1_n)=1$,
\item $\Phi(A)\geqslant0$ whenever $A$ is positive semi-definite.
\end{enumerate}
\end{definition}

If $\Phi$ is a state, it is in particular a (bounded) linear form (Remark~\ref{rem:finitedim}).
Since $M_n(\C)$ is moreover a Hilbert space by means of the Hilbert--Schmidt scalar product,
$\Phi$ has a special shape guaranteed by Riesz' Theorem~\ref{thm:Riesz}! This is the content
of the next proposition.

\begin{proposition}\label{prop:density}
Let $\Phi$ be a state on $M_n(\C)$. Then there exists a unique matrix $A\in M_n(\C)$ such
that $\Phi(B)=\tr(AB)$ for any $B\in M_n(\C)$. The matrix $A$ is normalized by $\tr(A)=1$
and is positive semi-definite.

Conversely, any such matrix $A$ defines a state $\Phi$ by the formula $\Phi(B)=\tr(AB)$. The
correspondence of normalized positive semi-definite matrices $A$ and states $\Phi$ is bijective.

$A$ is called the \textbf{\emph{density matrix}} of the state $\Phi$.
\end{proposition}

\begin{proof}
We first show the more difficult first part. Since $\Phi:M_n(\C)\to\C$ is a (bounded)
linear form on a Hilbert space, by Riesz' Theorem~\ref{thm:Riesz} there exists some
$A_0\in M_n(\C)$ such that $\Phi(B)=\ip{A_0}{B}=\tr(A_0^*B)$ for any $B\in M_n(\C)$. It
suffices to put $A=A_0^*$. We have $\tr(A)=\tr(A1_n)=\Phi(1_n)=1$. It remains to show that
$A$ is positive \emph{semi-}definite. Let $x\in\C^n$ be arbitrary. We associate the bounded
linear operator $T_x:\C^n\to\C^n$, $T_x(y)=\ip{x}{y}x$ (which is in fact equal to
$\ketbra{x}{x}$ from Example~\ref{ex:ketbra}). Then $T_x$ (or its associated matrix) is
positive semi-definite. Indeed, for any $y\in\C^n$, we have
\[
\ip{y}{T_x(y)}=\ip{y}{\ip{x}{y}x}=\ip{x}{y}\ip{y}{x}=|\ip{x}{y}|^2\ge0 .
\]
Now, the matrix of $T_x$ is equal to $(x_i\conj{x_j})_{ij}$ (compare Exercises 1.24 and 4.2):
indeed $(T_x)_{ij}=\ip{e_i}{T_xe_j}=\ip{x}{e_j}\,x_i=x_i\conj{x_j}$. Thus, with $A=(a_{ij})_{ij}$,
we have
\[
0\le\Phi(T_x)=\tr(AT_x)=\tr\Big(\Big(\sum_ja_{ij}x_j\conj{x_k}\Big)_{ik}\Big)
=\sum_i\sum_ja_{ij}x_j\conj{x_i}=\ip{x}{Ax}.
\]
Since $x$ was arbitrary, $A$ is positive semi-definite.

\emph{Uniqueness.} If $\tr(AB)=\tr(A'B)$ for all $B$, take $B=(A-A')^*$:
$\|(A-A')^*\|_{HS}^2=\tr((A-A')(A-A')^*)=0$, so $A=A'$.

\emph{Converse.} Let $A$ be positive semi-definite with $\tr A=1$ and $\Phi(B)=\tr(AB)$.
$\Phi$ is linear and $\Phi(1_n)=\tr A=1$. By Lemma~6.1a, $A=\sum_k\lambda_k T_{u_k}$ with
$\lambda_k\ge0$ and $(u_k)$ orthonormal. For any $u$ and $B$,
$\tr(T_uB)=\sum_{i,j}u_i\conj{u_j}b_{ji}=\ip{u}{Bu}$. Hence for $B$ positive semi-definite,
$\Phi(B)=\sum_k\lambda_k\ip{u_k}{Bu_k}\ge0$. So $\Phi$ is a state.

\emph{Bijectivity.} The map $A\mapsto\tr(A\,\cdot)$ from normalized positive semi-definite
matrices to states is well defined (converse), surjective (first part) and injective
(uniqueness).
\end{proof}

\begin{correction}[(proof of Prop.~6.3)]
(a) ``It remains to show that $A$ is positive \emph{definite}'' must read
``positive \emph{semi}-definite'': this is what the proposition claims and what the
argument gives; density matrices are in general not positive definite (see below).
(b) The original states that the matrix of $T_x$ is $(\conj{x_i}x_j)_{ij}$. This is wrong:
it is $(x_i\conj{x_j})_{ij}$, i.e.\ $xx^*$ (column times conjugated row), consistent with
Exercise~4.2 and with Example~8.4 where the original itself uses
$(\phi(t)_i\conj{\phi(s)_j})_{ij}$ for $\ketbra{\phi(t)}{\phi(s)}$. With the wrong matrix the
original obtains $\tr(AT_x)=\sum_{i,j}a_{ij}\conj{x_j}x_i$, which is $\ip{\conj x}{A\conj x}$
rather than $\ip{x}{Ax}$ (the conclusion survives only because $\conj x$ also runs through
all of $\C^n$; this is presumably the reason for the ``(!)'' in ``We infer(!)'').
(c) The uniqueness and the converse were not proved; see above.
\end{correction}

The above proof contains a positive semi-definite operator $T_x$. If $\|x\|_{\C^n}=1$, we
have $\tr(T_x)=\sum_i|x_i|^2=1$, and thus, by Proposition~\ref{prop:density}, $T_x$ defines a
state on $M_n(\C)$, $\Phi(B)=\tr(T_xB)=\ip{x}{Bx}$. Such a state is called \textbf{\emph{pure}}.
Note that $T_x$ has rank one, hence for $n\ge2$ it is \emph{not} positive definite.

\begin{exercise}\label{exo:mixed}
More generally, let $x_1,\dots,x_N\in\C^n$ be normalized vectors. Let
$\lambda_1,\dots,\lambda_N\geqslant0$ such that $\sum_{k=1}^N\lambda_k=1$. Check that
$\sum_{k=1}^N\lambda_kT_{x_k}$ defines a density matrix. If the $x_k$ span an at least
two-dimensional subspace and all $\lambda_k>0$, the corresponding state is called
\textbf{\emph{mixed}}.
\end{exercise}
\begin{correction}[(Ex.~6.4)]
The original upper summation limit is $n$ (the dimension) instead of $N$ (the number of vectors).
\end{correction}
\begin{solution}
Let $\rho=\sum_{k=1}^N\lambda_kT_{x_k}$. It is positive semi-definite:
$\ip{y}{\rho y}=\sum_k\lambda_k|\ip{x_k}{y}|^2\ge0$. And
$\tr\rho=\sum_k\lambda_k\tr T_{x_k}=\sum_k\lambda_k\|x_k\|^2=\sum_k\lambda_k=1$. By
Proposition~\ref{prop:density} it is a density matrix.

\emph{Why ``mixed'' is not pure.} Let all $\lambda_k>0$. If $\rho y=0$, then
$0=\ip{y}{\rho y}=\sum_k\lambda_k|\ip{x_k}{y}|^2$, so $y\perp x_k$ for all $k$; conversely,
if $y\perp x_k$ for all $k$, then $\rho y=\sum_k\lambda_k\ip{x_k}{y}x_k=0$. Hence
$\Ker\rho=\spn\{x_k\}^\perp$ and $\rank\rho=\dim\spn\{x_1,\dots,x_N\}\ge2$, whereas a pure
density matrix $T_x$ has rank $1$. Since the density matrix of a state is unique, a mixed
state is not pure.

\emph{Useful criterion (purity).} For a density matrix $\rho$ with eigenvalues
$p_k\ge0$, $\sum p_k=1$: $\tr(\rho^2)=\sum_kp_k^2\le1$, with equality iff one $p_k=1$ and
the others vanish, i.e.\ iff $\rho$ is pure.
\end{solution}

\begin{exercise}\label{exo:65}
Let
\[
A=\begin{pmatrix}2&1&-1\\1&5&1\\-1&1&1\end{pmatrix},\qquad x=(1,2,0),\quad y=(1,-1,-1).
\]
\begin{enumerate}
\item Show that $A=T_x+T_y$, where $T_x,T_y$ refer to the above notation, that is,
$T_x:\C^3\to\C^3$, $T_x(z)=\ip{x}{z}x$.
\item Deduce that $A$ is positive \emph{semi}-definite (but not positive definite) and that
$\Phi:M_3(\C)\to\C$, $B\mapsto\frac18\tr(AB)$ defines a mixed state.
\item Calculate
\[
\Phi\left(\begin{pmatrix}3&2&0\\0&-1&-1\\1&0&0\end{pmatrix}\right).
\]
\end{enumerate}
\end{exercise}

\begin{correction}[(Ex.~6.5) -- three errors]
(a) The original matrix is
$A_{\rm orig}=\begin{psmallmatrix}2&-1&-1\\-1&5&1\\-1&1&1\end{psmallmatrix}$. But
$T_x=xx^T=\begin{psmallmatrix}1&2&0\\2&4&0\\0&0&0\end{psmallmatrix}$ and
$T_y=yy^T=\begin{psmallmatrix}1&-1&-1\\-1&1&1\\-1&1&1\end{psmallmatrix}$, so
$T_x+T_y=\begin{psmallmatrix}2&1&-1\\1&5&1\\-1&1&1\end{psmallmatrix}$: the entries
$a_{12}=a_{21}$ have the wrong sign. (One checks that no choice $x=(x_1,x_2,0)$, $y\in\R^3$
gives $A_{\rm orig}$: the conditions $y_3^2=1$, $y_1y_3=-1$, $y_2y_3=1$, $x_1^2+y_1^2=2$,
$x_2^2+y_2^2=5$ force $|x_1|=1,|x_2|=2$, while $a_{12}=x_1x_2+y_1y_2=-1$ forces $x_1x_2=0$.)

(b) ``Deduce that $A$ is positive definite'' is false for $A=T_x+T_y$: a sum of two rank-one
matrices in $\C^3$ has rank $\le2$, so it has a non-trivial kernel. Here
$\det A=0$ and $A(-2,1,-3)^T=0$. $A$ is positive \emph{semi}-definite.
(Curiously, the misprinted $A_{\rm orig}$ \emph{is} positive definite: its leading minors
are $2,9,4>0$.)

(c) ``$\Phi:M_3(\C)\to M_3(\C)$, $B\mapsto\tr(AB)$'': a state takes values in $\C$, and
$\tr A=8\ne1$, so $\tr(A\,\cdot)$ is not normalized. The state is
$\Phi(B)=\frac18\tr(AB)$, with density matrix $\rho=\frac18A$.
\end{correction}

\begin{solution}
1. Since $x,y$ are real, $T_x$ has matrix $(x_ix_j)_{ij}$ and $T_y$ has matrix $(y_iy_j)_{ij}$,
computed in the correction above; their sum is $A$.

2. $T_x,T_y$ are positive semi-definite (proof of Prop.~\ref{prop:density}), hence so is
$A$: $\ip{z}{Az}=|\ip{x}{z}|^2+|\ip{y}{z}|^2\ge0$. It is not positive definite: for
$z=x\times y=(-2,1,-3)$, which is orthogonal to $x$ and $y$, $Az=0$. Its eigenvalues are
$0$ and $4\pm\sqrt2$. Now $\|x\|^2=5$, $\|y\|^2=3$, $\tr A=8$, and with
$\hat x=x/\sqrt5$, $\hat y=y/\sqrt3$,
\[
\rho=\tfrac18A=\tfrac58\,T_{\hat x}+\tfrac38\,T_{\hat y},
\]
which is of the form of Exercise~\ref{exo:mixed} with $\lambda_1=\frac58$,
$\lambda_2=\frac38>0$ and $\hat x,\hat y$ linearly independent. Hence
$\Phi(B)=\tr(\rho B)$ is a \emph{mixed} state (and $\rank\rho=2$).

3. Using $\tr(T_xM)=\ip{x}{Mx}$ (proof of Prop.~\ref{prop:density}), with $M$ the given
matrix: $Mx=(7,-2,1)$, $\ip{x}{Mx}=7-4+0=3$; $My=(1,2,1)$, $\ip{y}{My}=1-2-1=-2$. Thus
\[
\Phi(M)=\tfrac18\big(\ip{x}{Mx}+\ip{y}{My}\big)=\tfrac18(3-2)=\tfrac18 .
\]
(For comparison: with the misprinted $A_{\rm orig}$, which is also a valid --- full rank,
hence mixed --- density matrix after division by its trace $8$, one would get
$\frac18\tr(A_{\rm orig}M)=-\frac38$. Without the normalization one would get
$\tr(AM)=1$.)
\end{solution}

\begin{proposition}\label{prop:bloch}
The density matrices on $M_2(\C)$ are exactly the matrices of the form
$A=\frac12(1_2+a_x\sigma_x+a_y\sigma_y+a_z\sigma_z)$ with $a_x,a_y,a_z\in\R$ and
$\|a\|\le1$, where $a=(a_x,a_y,a_z)$. The pure states correspond to vectors $a\in\R^3$ such
that $\|a\|=1$, the mixed states to vectors such that $\|a\|<1$.
$\{a\in\R^3:\|a\|\le1\}$ is called the \emph{Bloch ball}, $\{a\in\R^3:\|a\|=1\}$ is called
the \emph{Bloch sphere}.
\end{proposition}

\begin{addition}[(proof of Proposition~6.6)]
Let $A$ be a density matrix. By Exercise~4.17, $A=\frac12(c_01+a_x\sigma_x+a_y\sigma_y+a_z\sigma_z)$
for some $c_0,a_j\in\C$, with $c_0=\ip{1}{A}=\tr A=1$ and $a_j=\ip{\sigma_j}{A}=\tr(\sigma_jA)$.
Since $A$ is self-adjoint (Lemma~6.1a) and the $\sigma_j$ are self-adjoint,
$\conj{a_j}=\conj{\tr(\sigma_jA)}=\tr((\sigma_jA)^*)=\tr(A\sigma_j)=a_j$, so $a_j\in\R$.
Conversely, every such matrix is self-adjoint with trace $1$. It remains to see when it is
positive semi-definite. Put $a\cdot\sigma=a_x\sigma_x+a_y\sigma_y+a_z\sigma_z$. Using
$\sigma_j^2=1$ and $\sigma_j\sigma_k+\sigma_k\sigma_j=0$ for $j\ne k$ (from the products in
Exercise~4.17), $(a\cdot\sigma)^2=\|a\|^21$. So the eigenvalues of $a\cdot\sigma$ are in
$\{\pm\|a\|\}$, and since $\tr(a\cdot\sigma)=0$, they are $+\|a\|$ and $-\|a\|$. Hence
the eigenvalues of $A$ are $\frac{1\pm\|a\|}2$, and $A\ge0$ iff $\|a\|\le1$ (Lemma~6.1a).

$A$ is pure iff $A=T_x$ for a unit vector $x$ iff $\rank A=1$ (a positive semi-definite
rank-one matrix with trace $1$ is $\lambda T_x$ with $\lambda=1$) iff one eigenvalue
vanishes iff $\|a\|=1$. Otherwise ($\|a\|<1$) both eigenvalues are $>0$, $A$ has rank $2$
and the state is mixed. (Equivalently $\tr A^2=\frac{1+\|a\|^2}2$.)

\emph{Pure states as points on the sphere.} Every unit vector of $\C^2$ is, up to a phase,
$x=(\cos\frac\theta2,\,e^{i\varphi}\sin\frac\theta2)$ with $\theta\in[0,\pi]$,
$\varphi\in[0,2\pi)$, and a direct computation gives $T_x=\frac12(1+a\cdot\sigma)$ with
$a=(\sin\theta\cos\varphi,\sin\theta\sin\varphi,\cos\theta)$. For instance
$\ket0=(1,0)\leftrightarrow a=(0,0,1)$ (north pole), $\ket1\leftrightarrow(0,0,-1)$, and the
maximally mixed state $\frac121_2\leftrightarrow a=0$ (centre).
\end{addition}
%======================================================================
\section{Metric spaces}

\begin{definition}
Let $X$ be a non-empty set and $d:X\times X\to[0,\infty)$ be a function. Then $d$ is
called a \emph{metric} (or distance) on $X$ and $(X,d)$ (or simply $X$) is called a
\emph{metric space} provided that
\begin{enumerate}
\item $d(x,y)=0\iff x=y$ (definiteness),
\item $d(x,y)=d(y,x)$ for all $x,y\in X$ (symmetry),
\item $d(x,z)\le d(x,y)+d(y,z)$ for all $x,y,z\in X$ (triangle inequality).
\end{enumerate}
\end{definition}

\begin{definition}
Let $(X,d)$ be a metric space. Let $r>0$ and $x_0\in X$.
\begin{enumerate}
\item The set $B(x_0,r)=\{y\in X:\ d(x_0,y)<r\}$ is called \emph{open ball} in $X$ with
center $x_0$ and radius $r$.
\item The set $\conj B(x_0,r)=\{y\in X:\ d(x_0,y)\le r\}$ is called \emph{closed ball} in
$X$ with center $x_0$ and radius $r$.
\end{enumerate}
\end{definition}
\begin{correction}[(Def.~7.2)]
The original allows $r\ge0$; for $r=0$ the open ball is empty and the closed ball is
$\{x_0\}$. This is harmless but unusual: balls are normally taken with $r>0$.
\end{correction}

\begin{exercise}
Let $(X,d)$ be a metric space. Let $x,y,z\in X$. Show the reverse triangle inequality
\[
|d(x,y)-d(y,z)|\le d(x,z).
\]
\end{exercise}
\begin{solution}
By the triangle inequality and symmetry, $d(x,y)\le d(x,z)+d(z,y)=d(x,z)+d(y,z)$, i.e.\
$d(x,y)-d(y,z)\le d(x,z)$. Exchanging the roles of $x$ and $z$:
$d(z,y)\le d(z,x)+d(x,y)$, i.e.\ $d(y,z)-d(x,y)\le d(x,z)$. Both together give
$|d(x,y)-d(y,z)|\le d(x,z)$.
\end{solution}

\begin{example}
Let $X$ be a Hilbert space (or more generally a normed space). Then $d(x,y)=\|x-y\|$ defines
a metric on $X$. With this example in mind, it is instructive to note that the concept of a
metric space is a generalization of the concept of a normed space (in particular of a
Hilbert space).
\end{example}
\begin{proof}
$d(x,y)=0\iff x-y=0$; $d(y,x)=\|-(x-y)\|=|-1|\,\|x-y\|=d(x,y)$;
$\|x-z\|=\|(x-y)+(y-z)\|\le\|x-y\|+\|y-z\|$.
\end{proof}

\begin{example}\label{ex:sncf}
Let $X=\R^2$ with Euclidean norm $|\cdot|$. Define
\[
d(x,y)=\begin{cases}|x-y|&:\ x,y \text{ are collinear with } 0\\ |x|+|y|&:\ \text{otherwise}\end{cases}.
\]
Then $d$ defines a metric on $X$, the \emph{SNCF metric} (or ``French railway metric'':
to travel between two cities not on the same line through Paris $=0$, one has to change
trains in Paris). Open balls $B(0,r)$ are the usual open discs in $\R^2$. Open balls
$B(x_0,r)$ for $0<r\le|x_0|$ are open segments lying on the line joining $x_0$ to $0$.
\end{example}

\begin{correction}[(Ex.~7.5)]
``$x,y$ are colinear'' must be understood as ``$x$, $y$ and $0$ lie on a common line'', i.e.\
$x$ and $y$ are linearly dependent (any two points are collinear in the usual sense!).
``Open circles'' should be ``open discs''. The statement on balls also holds for
$r=|x_0|$. The metric axioms were not verified; see below.
\end{correction}

\begin{addition}[(proofs for Example~7.5)]
Note first that always $|x-y|\le d(x,y)\le|x|+|y|$.
\emph{Definiteness:} if $x,y$ are collinear with $0$, $d(x,y)=|x-y|$; otherwise $x,y\ne0$
(the vector $0$ is collinear with everything) and $d(x,y)=|x|+|y|>0$.
\emph{Symmetry} is clear.
\emph{Triangle inequality} $d(x,z)\le d(x,y)+d(y,z)$:
\begin{itemize}
\item If $x,z$ are collinear with $0$: $d(x,z)=|x-z|\le|x-y|+|y-z|\le d(x,y)+d(y,z)$.
\item Otherwise $d(x,z)=|x|+|z|$ and $x,z\ne0$. If $y=0$:
$d(x,0)+d(0,z)=|x|+|z|$. If $y\ne0$, $y$ cannot be collinear with both $x$ and $z$ (they
would both lie on the line $\R y$). Say $y$ is not collinear with $x$: then
$d(x,y)+d(y,z)\ge|x|+|y|+|y-z|\ge|x|+|y|+(|z|-|y|)=|x|+|z|$. The case where $y$ is not
collinear with $z$ is symmetric.
\end{itemize}
\emph{Balls.} $B(0,r)=\{y:|y|<r\}$ since $0$ is collinear with every $y$. For $x_0\ne0$
and $y$ not on the line $\R x_0$: $d(x_0,y)=|x_0|+|y|\ge|x_0|$. Hence for $r\le|x_0|$,
$B(x_0,r)=\{y\in\R x_0:\ |y-x_0|<r\}$, an open segment of length $2r$ centred at $x_0$
on the line through $0$ and $x_0$ (not containing $0$). For $r>|x_0|$ one obtains this
segment together with the Euclidean disc $\{y:|y|<r-|x_0|\}$.
(The metric is thus very different from the Euclidean one, although it induces the
same distances along lines through $0$.)
\end{addition}

%======================================================================
\section{Reproducing kernels}

\begin{definition}\label{def:psdkernel}
Let $X$ be a non-empty set. Let $k:X\times X\to\C$ be a function. Then $k$ is called
\emph{positive semi-definite} if for all $n\in\N$, $x_1,\dots,x_n\in X$ and
$c_1,\dots,c_n\in\C$,
\[
\sum_{i,j=1}^nc_i\conj{c_j}\,k(x_i,x_j)\ge0 .
\]
Moreover, $k$ is called \emph{positive definite} if the above sum is always strictly
positive, provided that the $x_i$ are pairwise distinct and not all $c_1,\dots,c_n$ are
equal to $0$.
\end{definition}

\begin{correction}[(Def.~8.1)]
For positive \emph{definiteness} one must require the points $x_1,\dots,x_n$ to be pairwise
distinct; otherwise no kernel could be positive definite (take $x_1=x_2$, $c_1=1$,
$c_2=-1$: the sum is $0$).
\end{correction}

\begin{addedthm}{Lemma 8.1a}
Let $k$ be positive semi-definite. Then for all $x,y\in X$:
$k(x,x)\ge0$, $k(y,x)=\conj{k(x,y)}$ (Hermitian symmetry), and
$|k(x,y)|^2\le k(x,x)k(y,y)$.
\end{addedthm}
\begin{proof}
$n=1$ gives $|c|^2k(x,x)\ge0$. With $n=2$, $x_1=x$, $x_2=y$:
$q=|c_1|^2k(x,x)+|c_2|^2k(y,y)+c_1\conj{c_2}k(x,y)+c_2\conj{c_1}k(y,x)\ge0$, in particular
real. With $c_1=c_2=1$: $k(x,y)+k(y,x)\in\R$. With $c_1=1$, $c_2=i$:
$-ik(x,y)+ik(y,x)\in\R$, i.e.\ $k(y,x)-k(x,y)\in i\R$. The first says
$\Imag k(y,x)=-\Imag k(x,y)$, the second $\Real k(y,x)=\Real k(x,y)$; so
$k(y,x)=\conj{k(x,y)}$. Finally, the $2\times2$ matrix
$\begin{psmallmatrix}k(x,x)&k(y,x)\\k(x,y)&k(y,y)\end{psmallmatrix}$ is positive
semi-definite (this is $q\ge0$), so its determinant (product of its eigenvalues $\ge0$,
Lemma~6.1a) is $k(x,x)k(y,y)-|k(x,y)|^2\ge0$.
\end{proof}

\begin{example}\label{ex:feature}
Let $X$ be a non-empty set, $H$ be a Hilbert space and $A:X\to H$ be any mapping. Then
$k(x,y)=\ip{A(y)}{A(x)}$ defines a positive semi-definite function. Indeed, if
$x_1,\dots,x_n\in X$ and $c_1,\dots,c_n\in\C$, we have
\begin{align*}
\sum_{i,j=1}^nc_i\conj{c_j}k(x_i,x_j)&=\sum_{i,j=1}^nc_i\conj{c_j}\ip{A(x_j)}{A(x_i)}
=\Big\langle\sum_{j=1}^nc_jA(x_j),\sum_{i=1}^nc_iA(x_i)\Big\rangle\\
&=\Big\|\sum_{i=1}^nc_iA(x_i)\Big\|^2\ge0 .
\end{align*}
\end{example}

In fact, we will see a converse to the above example in the final theorem of this lecture
(Theorem~8.7 below). Let us turn to an example of a positive semi-definite function arising
in quantum mechanics: the quantum kernel.

\begin{definition}\label{def:qkernel}
Let $X$ be a set (of configurations), $\C^n$ be a finite-dimensional Hilbert space and
$\phi:X\to\C^n$ be a representation of the configurations, so that we must have that
$\|\phi(x)\|_{\C^n}=1$ for all $x\in X$. Then the associated \emph{data-encoding feature map} is
\[
\rho:X\to M_n(\C),\qquad\rho(x)=\ketbra{\phi(x)}{\phi(x)}
\]
(see Example~\ref{ex:ketbra} and the proof of Proposition~\ref{prop:density}). Thus, $\rho$
takes values in the pure states of $M_n(\C)$. Moreover,
\[
k:X\times X\to\C,\qquad k(x,y)=\tr(\rho(y)\rho(x))\overset{!}{=}\ip{\rho(y)}{\rho(x)}_{\text{Hilbert--Schmidt}}
\]
is called the \emph{quantum kernel}. It is a positive semi-definite function according to
Example~\ref{ex:feature}.
\end{definition}

\begin{addition}[(the ``$\overset{!}{=}$'' and a closed formula)]
$\ip{\rho(y)}{\rho(x)}_{HS}=\tr(\rho(y)^*\rho(x))=\tr(\rho(y)\rho(x))$ because
$\rho(y)=T_{\phi(y)}$ is self-adjoint. Example~\ref{ex:feature} applies with $H=M_n(\C)$,
$A=\rho$. Moreover, using $\ketbra uv\,\ketbra wz=\ip vw\,\ketbra uz$ and
$\tr\ketbra uv=\sum_iu_i\conj{v_i}=\ip vu$, we obtain for every quantum kernel
\[
k(x,y)=\ip{\phi(y)}{\phi(x)}\tr\ketbra{\phi(y)}{\phi(x)}=\ip{\phi(y)}{\phi(x)}\ip{\phi(x)}{\phi(y)}
=|\ip{\phi(x)}{\phi(y)}|^2\in[0,1],
\]
the \emph{fidelity} between the encoded states; it does not depend on global phases of
$\phi(x)$, as it should be physically. This is exactly what a quantum computer estimates
(e.g.\ by the swap test or by measuring $U(y)^*U(x)\ket0$ and counting the outcome $\ket0$).
\end{addition}

\begin{example}\label{ex:qkernel}
Let us calculate the quantum kernel in an example. Let $n=2$ and for $t\in\R$
\[
U(t)=e^{-i\frac t2\sigma_x}=\begin{pmatrix}\cosh(-i\frac t2)&\sinh(-i\frac t2)\\\sinh(-i\frac t2)&\cosh(-i\frac t2)\end{pmatrix}
=\begin{pmatrix}\cos\frac t2&-i\sin\frac t2\\-i\sin\frac t2&\cos\frac t2\end{pmatrix}
\]
be the matrix exponential from Exercise~\ref{exo:expsx}. The columns of $U(t)$ are
normalized and orthogonal to each other, so $U(t)$ is indeed unitary according to
Exercise~\ref{exo:unitcols}. Let then $\phi:\R\to\C^2$,
$t\mapsto U(t)(1,0)=\big(\cos\frac t2,-i\sin\frac t2\big)$. As $U(t)$ is unitary,
$\|\phi(t)\|=1$ $(t\in\R)$, so that $\phi$ is a valid representation of configurations.
Moreover, the associated quantum kernel is
\begin{align*}
k(s,t)&=\tr(\rho(t)\rho(s))=\tr\big(\ketbra{\phi(t)}{\phi(t)}\cdot\ketbra{\phi(s)}{\phi(s)}\big)
=\tr\big(\ket{\phi(t)}\ip{\phi(t)}{\phi(s)}\bra{\phi(s)}\big)\\
&=\ip{\phi(t)}{\phi(s)}\tr\big((\phi(t)_i\conj{\phi(s)_j})_{ij}\big)
=\ip{\phi(t)}{\phi(s)}\big(\phi(t)_1\conj{\phi(s)_1}+\phi(t)_2\conj{\phi(s)_2}\big)\\
&=|\ip{\phi(t)}{\phi(s)}|^2=\Big(\cos\frac t2\cos\frac s2+\sin\frac t2\sin\frac s2\Big)^2
=\Big(\cos\frac{t-s}2\Big)^2 .
\end{align*}
We deduce from Definition~\ref{def:qkernel} that the function
$k(s,t)=\big(\cos\frac{t-s}2\big)^2$ is positive semi-definite $\R\times\R\to\C$.
\end{example}

\begin{addition}[(details and remarks on Example~8.4)]
\emph{Step $\ip{\phi(t)}{\phi(s)}$:}
$\conj{\cos\frac t2}\cos\frac s2+\conj{(-i\sin\frac t2)}(-i\sin\frac s2)
=\cos\frac t2\cos\frac s2+(i)(-i)\sin\frac t2\sin\frac s2=\cos\frac{t-s}2$.
\emph{Physical meaning:} $\phi(t)$ is the qubit $\ket0$ rotated by the angle $t$ about the
$x$-axis of the Bloch sphere; its Bloch vector is $(0,-\sin t,\cos t)$.
\emph{Direct check of positivity:} $k(s,t)=\frac{1+\cos(t-s)}2
=\frac12+\frac12\cos t\cos s+\frac12\sin t\sin s$ is of the form of Example~\ref{ex:feature}
with $A(t)=\frac1{\sqrt2}(1,\cos t,\sin t)\in\R^3$.
\emph{Not positive definite:} $k(s,t)=1$ if $t-s\in2\pi\Z$, so for $x_1=0$, $x_2=2\pi$,
$c=(1,-1)$ the sum in Def.~8.1 vanishes. (Indeed $\phi(2\pi)=-\phi(0)$: same state.)
\end{addition}

Quantum kernels, or more generally positive semi-definite kernels, can be used to
construct special Hilbert spaces, the reproducing kernel Hilbert spaces (RKHS).

\begin{definition}
Let $X$ be a non-empty set. A \emph{reproducing kernel Hilbert space (RKHS)} over $X$ is by
definition a Hilbert space consisting of functions over $X$ (with the usual pointwise
addition and scalar multiplication) such that for any $x\in X$, the point evaluation
$H\to\C$, $f\mapsto f(x)$ is a bounded form (it is always linear).
\end{definition}

\begin{example}\label{ex:rkhs}
Let $k:X\times X\to\C$ be a positive semi-definite kernel. Let $H_0$ be the linear subspace
of all the functions from $X$ to $\C$ which is generated by the functions
$\{k_x:\ x\in X\}$, where $k_x(y)=k(x,y)$. If $f=\sum_jc_jk_{x_j}$ and $g=\sum_ld_lk_{y_l}$
belong to $H_0$ (finite sums), then
\begin{equation}\label{eq:81}
\sum_l\conj{d_l}f(y_l)=\sum_{j,l}c_j\conj{d_l}\,k(x_j,y_l)=\sum_jc_j\conj{g(x_j)},
\end{equation}
so that this quantity does not depend on the chosen representations of $f,g$ (which may be
not unique). Indeed, the left-hand side depends only on $f$ (and the chosen representation of
$g$), the right-hand side only on $g$ (and the representation of $f$). We denote this
quantity by $\ip{g}{f}$. Then
\[
\ip ff=\sum_{j,l}c_j\conj{c_l}\,k(x_j,x_l)\ge0
\]
by the assumption that $k$ is positive semi-definite. It is clear from \eqref{eq:81} that the
form $\ip\cdot\cdot$ is linear in its second argument and antilinear in its first. An
immediate consequence of \eqref{eq:81} (take only one index in $l$, $d_1=1$ and $y_1=x$) is
the \emph{reproducing property}
\[
\ip{k_x}{f}=f(x)\qquad(f\in H_0,\ x\in X).
\]
This implies
\begin{equation}\label{eq:82}
\ip{k_x}{k_y}=k(y,x).
\end{equation}
Moreover, in a similar manner to Proposition~\ref{prop:CS} (precisely: by Lemma~1.7a), one can show
\[
|f(x)|^2=|\ip{k_x}{f}|^2\le\ip ff\cdot k(x,x),
\]
so that $\ip ff=0$ if and only if $f$ is identically zero. Therefore, $\ip\cdot\cdot$ is a
scalar product. One can complete $H_0$ (i.e.\ enlarge $H_0$ by limits of Cauchy sequences in
$H_0$) to obtain a Hilbert space $H\supseteq H_0$ with the same scalar product. Moreover,
$H$ is again a linear subspace of the functions from $X$ to $\C$.

To sum up, there is a Hilbert space $H\subseteq\C^X$ and a mapping $x\mapsto k_x$ from $X$ to
$H$ such that $k(y,x)=\ip{k_x}{k_y}$ for $x,y\in X$ according to \eqref{eq:82}. Thus, every
positive semi-definite kernel $k$ arises in the form of Example~\ref{ex:feature} (with
$A(x)=k_x$)! Finally, $H$ is a RKHS. Indeed, the point evaluation in any $x\in X$ is given by
$f(x)=\ip{k_x}{f}$, which is continuous by (the easy direction of) Riesz' Theorem~\ref{thm:Riesz}.
\end{example}

\begin{correction}[(Ex.~8.6) -- gaps in the argument]
(a) The second equality in \eqref{eq:81},
$\sum_{j,l}c_j\conj{d_l}k(x_j,y_l)=\sum_jc_j\conj{g(x_j)}$, needs
$\conj{g(x_j)}=\sum_l\conj{d_l}\,\conj{k(y_l,x_j)}=\sum_l\conj{d_l}\,k(x_j,y_l)$, i.e.\ the
Hermitian symmetry $\conj{k(y,x)}=k(x,y)$. This is true for every positive semi-definite
kernel, but has to be proved (Lemma~8.1a). The same symmetry shows
$\ip gf=\conj{\ip fg}$.
(b) Cauchy--Schwarz (Prop.~1.7) was proved for scalar products, i.e.\ assuming
definiteness, which is precisely what one wants to deduce here. One needs the version for
semi-definite forms (Lemma~1.7a), applied with $\ip{k_x}{k_x}=k(x,x)$.
(c) The facts that the abstract completion of $H_0$ can be realized as a space of
\emph{functions} on $X$ (with the reproducing property extended) is the non-trivial part
of the Moore--Aronszajn theorem; it is proved below.
\end{correction}

\begin{addedthm}{Theorem 8.7 (Moore--Aronszajn)}
Let $k:X\times X\to\C$ be positive semi-definite. Then there is a unique Hilbert space $H$ of
functions $X\to\C$ such that $k_x\in H$ for all $x\in X$ and $f(x)=\ip{k_x}{f}_H$ for all
$f\in H$, $x\in X$. $H$ is a RKHS, $H_0$ is dense in $H$, and $k(x,y)=\ip{k_y}{k_x}_H$, i.e.\
$k$ is of the form of Example~\ref{ex:feature} with feature map $A(x)=k_x$.
\end{addedthm}
\begin{proof}
Let $H_0$ with $\ip\cdot\cdot$ be as in Example~\ref{ex:rkhs} (a scalar product space).

\emph{Step 1: Cauchy sequences converge pointwise.} If $(f_n)$ is Cauchy in $H_0$, then
$|f_n(x)-f_m(x)|=|\ip{k_x}{f_n-f_m}|\le\sqrt{k(x,x)}\,\|f_n-f_m\|$, so $f(x):=\lim_nf_n(x)$
exists for every $x$. Let $H$ be the set of all functions $f:X\to\C$ obtained in this way.
$H$ is a vector space containing $H_0$ (constant sequences).

\emph{Step 2: a Cauchy sequence tending pointwise to $0$ tends to $0$ in norm.} Let $(g_n)$
be Cauchy in $H_0$ with $g_n(x)\to0$ for all $x$. It is bounded, $\|g_n\|\le M$. Given
$\epsilon>0$ choose $N$ with $\|g_n-g_N\|\le\epsilon$ for $n\ge N$, and write
$g_N=\sum_jc_jk_{x_j}$ (finite sum). Then for $n\ge N$,
\[
\|g_n\|^2=\ip{g_n-g_N}{g_n}+\ip{g_N}{g_n}\le\epsilon M+\Big|\sum_j\conj{c_j}\,g_n(x_j)\Big|
\xrightarrow[n\to\infty]{}\epsilon M .
\]
Hence $\limsup\|g_n\|^2\le\epsilon M$ for every $\epsilon>0$, i.e.\ $\|g_n\|\to0$.

\emph{Step 3: scalar product on $H$.} For $f,g\in H$ with approximating Cauchy sequences
$(f_n),(g_n)$ in $H_0$, the limit $\ip fg_H:=\lim_n\ip{f_n}{g_n}$ exists, since
$|\ip{f_n}{g_n}-\ip{f_m}{g_m}|\le\|f_n-f_m\|\|g_n\|+\|f_m\|\|g_n-g_m\|$. It does not depend on
the chosen sequences: if $(f'_n)$ also converges pointwise to $f$, then $(f_n-f'_n)$ is Cauchy
and tends pointwise to $0$, hence $\|f_n-f'_n\|\to0$ by Step 2. Sesquilinearity, conjugate
symmetry and positivity pass to the limit. \emph{Reproducing property:}
$\ip{k_x}f_H=\lim_n\ip{k_x}{f_n}=\lim_nf_n(x)=f(x)$. \emph{Definiteness:} if $\ip ff_H=0$, then
$|f(x)|^2=|\ip{k_x}f_H|^2\le k(x,x)\ip ff_H=0$ (Cauchy--Schwarz in $H$), so $f=0$.

\emph{Step 4: $H_0$ is dense and $H$ is complete.} If $f_n\to f$ pointwise as above, then
$f-f_m$ is approximated by the Cauchy sequence $(f_n-f_m)_n$, so
$\|f-f_m\|_H^2=\lim_n\|f_n-f_m\|^2$, which is $\le\epsilon^2$ for $m$ large: $f_m\to f$ in
$H$. Now let $(h_m)$ be Cauchy in $H$; pick $g_m\in H_0$ with $\|h_m-g_m\|_H\le\frac1m$.
Then $(g_m)$ is Cauchy in $H_0$, so it converges pointwise to some $g\in H$, and
$g_m\to g$ in $H$; hence $h_m\to g$ in $H$.

\emph{Step 5: RKHS and uniqueness.} Point evaluations $f\mapsto\ip{k_x}f_H$ are bounded
(Exercise~1.20), so $H$ is a RKHS; and $\ip{k_y}{k_x}_H=k_x(y)=k(x,y)$. If $H'$ is another Hilbert space of functions with the
stated properties, then $H_0\subseteq H'$ with the same scalar product (by the reproducing
property, $\ip{k_x}{k_y}_{H'}=k(y,x)$), hence $H\subseteq H'$ isometrically (norm limits in
$H'$ are pointwise limits, by the reproducing property). If $f\in H'$ is orthogonal to $H$,
then $f(x)=\ip{k_x}{f}_{H'}=0$ for all $x$; hence $H'=H$.
\end{proof}

\begin{addeddef}{Example 8.8 (the RKHS of the quantum kernel of Example~8.4)}
For $k(s,t)=\cos^2\frac{t-s}2=\frac12+\frac12\cos s\cos t+\frac12\sin s\sin t$ we have
$k_s=\frac12\cdot1+\frac{\cos s}2\cos+\frac{\sin s}2\sin$, so $H_0\subseteq V:=\spn\{1,\cos,\sin\}$.
Since $k_0=\frac12(1+\cos)$, $k_\pi=\frac12(1-\cos)$, $k_{\pi/2}=\frac12(1+\sin)$, in fact
$H_0=V$, which is $3$-dimensional, hence complete: $H=V$. On $V$ define
$\ip{\alpha+\beta\cos+\gamma\sin}{\alpha'+\beta'\cos+\gamma'\sin}
:=2(\conj\alpha\alpha'+\conj\beta\beta'+\conj\gamma\gamma')$. Then
$\ip{k_s}{f}=2\big(\frac12\alpha+\frac{\cos s}2\beta+\frac{\sin s}2\gamma\big)=f(s)$ for
$f=\alpha+\beta\cos+\gamma\sin$: this is the reproducing property, so by uniqueness in
Theorem~8.7 this is the RKHS norm: $\|\alpha+\beta\cos+\gamma\sin\|_H^2=2(|\alpha|^2+|\beta|^2+|\gamma|^2)$.
In machine-learning terms: a kernel method (SVM, kernel ridge regression) with this quantum
kernel can only learn functions of the form $\alpha+\beta\cos t+\gamma\sin t$ --- the
expressivity of a quantum kernel is governed by its (here 3-dimensional) RKHS.
\end{addeddef}

\end{document}
